\label{kap:priroda}
\mathquote{Mraky nejsou koule, hory nejsou kužely, pobřeží nejsou kružnice, kůra není hladká a blesk se nešíří po přímce.}{Benoît Mandelbrot, Fraktální geometrie přírody, 1982}

Příroda matematiku nestudovala, a přesto se jí řídí. Epidemie se šíří podle diferenciálních rovnic z kapitoly~\ref{kap:diferencialni}. Počty zvířat mohou kolísat chaoticky, přestože se řídí jednoduchým vzorcem. Semínka slunečnice se rovnají podle čísla, které jsme potkali u Fibonacciho posloupnosti v kapitole~\ref{kap:posloupnosti}. Harmonie v hudbě je otázkou zlomků a mocnin dvojky z kapitoly~\ref{kap:exponencialni}. A k tomu dvě věci, které přírodní nejsou, ale bez matematiky by nefungovaly: GPS v telefonu, které musí počítat s Einsteinovou relativitou, a vyhledávač, který řadí stránky pomocí vlastního vektoru matice z kapitoly~\ref{kap:linearni-algebra}.

\section{Epidemie: model SIR}

Rozdělme populaci na tři skupiny: \textbf{náchylní} ($S$, \textit{susceptible}), \textbf{nakažení} ($I$, \textit{infected}) a \textbf{zotavení} ($R$, \textit{recovered}, sem patří i zemřelí). Počítáme podíly, takže $S + I + R = 1$.

\begin{definitionbox}
  \textbf{Model SIR} (Kermack a McKendrick, 1927):
  \begin{equation*}
    \frac{\mathrm{d}S}{\mathrm{d}t} = -\beta SI,\qquad
    \frac{\mathrm{d}I}{\mathrm{d}t} = \beta SI - \gamma I,\qquad
    \frac{\mathrm{d}R}{\mathrm{d}t} = \gamma I.
  \end{equation*}
  Číslo $\beta$ je rychlost nakažlivých kontaktů, $\gamma$ rychlost uzdravování ($\frac{1}{\gamma}$ je průměrná doba nemoci). \textbf{Základní reprodukční číslo} je
  \begin{equation*}
    R_0 = \frac{\beta}{\gamma}.
  \end{equation*}
\end{definitionbox}

Odkud se rovnice berou? Nakažení se setkávají s náchylnými. Počet takových setkání za jednotku času je úměrný $S\cdot I$ (jako u chemické reakce dvou látek), a část z nich vede k nákaze. Tak vzniká člen $\beta SI$, který ubírá z $S$ a přidává do $I$. Nakažení se zároveň uzdravují rychlostí $\gamma I$. A $R_0$ je počet lidí, které jeden nakažený nakazí v plně náchylné populaci: za den nakazí $\beta$ lidí a nemocný je průměrně $\frac{1}{\gamma}$ dní.

\begin{center}
  \begin{tikzpicture}[x=0.1cm, y=4cm]
    \draw[gray!30] (0,0.25) -- (100,0.25) (0,0.5) -- (100,0.5) (0,0.75) -- (100,0.75) (0,1) -- (100,1);
    \draw[->] (0,0) -- (104,0) node[right, font=\small] {$t$ (dny)};
    \draw[->] (0,0) -- (0,1.1);
    \foreach \y in {0.25,0.5,0.75,1} \node[left, font=\scriptsize] at (0,\y) {$\y$};
    \foreach \x in {0,20,...,100} \node[below, font=\scriptsize] at (\x,0) {$\x$};
    \draw[very thick, gray] plot[smooth] coordinates {(0,0.990) (2,0.983) (4,0.972) (6,0.957) (8,0.935) (10,0.904) (12,0.863) (14,0.810) (16,0.744) (18,0.667) (20,0.584) (22,0.501) (24,0.423) (26,0.353) (28,0.294) (30,0.246) (32,0.208) (34,0.178) (36,0.154) (38,0.136) (40,0.121) (42,0.109) (44,0.100) (46,0.093) (48,0.087) (50,0.082) (52,0.078) (54,0.075) (56,0.072) (58,0.070) (60,0.068) (62,0.067) (64,0.066) (66,0.064) (68,0.064) (70,0.063) (72,0.062) (74,0.062) (76,0.061) (78,0.061) (80,0.061) (82,0.060) (84,0.060) (86,0.060) (88,0.060) (90,0.060) (92,0.059) (94,0.059) (96,0.059) (98,0.059) (100,0.059)};
    \draw[very thick, colorofwisdom] plot[smooth] coordinates {(0,0.010) (2,0.015) (4,0.022) (6,0.032) (8,0.046) (10,0.065) (12,0.091) (14,0.123) (16,0.161) (18,0.201) (20,0.240) (22,0.272) (24,0.294) (26,0.303) (28,0.301) (30,0.290) (32,0.272) (34,0.250) (36,0.226) (38,0.202) (40,0.178) (42,0.157) (44,0.136) (46,0.118) (48,0.102) (50,0.088) (52,0.076) (54,0.065) (56,0.056) (58,0.047) (60,0.040) (62,0.035) (64,0.029) (66,0.025) (68,0.021) (70,0.018) (72,0.015) (74,0.013) (76,0.011) (78,0.009) (80,0.008) (82,0.007) (84,0.006) (86,0.005) (88,0.004) (90,0.004) (92,0.003) (94,0.003) (96,0.002) (98,0.002) (100,0.002)};
    \draw[very thick, black, dashed] plot[smooth] coordinates {(0,0.000) (2,0.002) (4,0.006) (6,0.011) (8,0.019) (10,0.030) (12,0.046) (14,0.067) (16,0.095) (18,0.131) (20,0.176) (22,0.227) (24,0.284) (26,0.344) (28,0.404) (30,0.464) (32,0.520) (34,0.572) (36,0.620) (38,0.662) (40,0.701) (42,0.734) (44,0.763) (46,0.789) (48,0.811) (50,0.830) (52,0.846) (54,0.860) (56,0.872) (58,0.882) (60,0.891) (62,0.899) (64,0.905) (66,0.910) (68,0.915) (70,0.919) (72,0.922) (74,0.925) (76,0.928) (78,0.930) (80,0.931) (82,0.933) (84,0.934) (86,0.935) (88,0.936) (90,0.937) (92,0.938) (94,0.938) (96,0.939) (98,0.939) (100,0.939)};
    \draw[dotted] (26.6,0) -- (26.6,0.304);
    \draw[dotted] (0,0.3333) -- (26.6,0.3333);
    \node[font=\scriptsize, anchor=west] at (1,0.43) {vrchol při $S = \frac{1}{R_0}$};
    \node[font=\small, gray] at (80,0.12) {$S$};
    \node[font=\small, colorofwisdom] at (38,0.2) {$I$};
    \node[font=\small] at (80,0.88) {$R$};
  \end{tikzpicture}
\end{center}

Obrázek je numerické řešení pro $\beta = 0{,}3$, $\gamma = 0{,}1$ (nemoc trvá průměrně $10$ dní, $R_0 = 3$), na začátku je nakaženo $1\,\%$ populace. Epidemie nejdřív roste téměř exponenciálně, kolem $27$. dne vrcholí s asi $30\,\%$ současně nakažených a pak odezní. Nakonec projde nemocí asi $94\,\%$ lidí.

\subsection{Kolektivní imunita}

Počet nakažených roste, právě když $\frac{\mathrm{d}I}{\mathrm{d}t} = I(\beta S - \gamma) > 0$, tedy když
\begin{equation*}
  S > \frac{\gamma}{\beta} = \frac{1}{R_0}.
\end{equation*}
Epidemie tedy nemůže propuknout, pokud je náchylných méně než $\frac{1}{R_0}$. Jinými slovy, imunní musí být aspoň podíl
\begin{equation*}
  p_c = 1 - \frac{1}{R_0}
\end{equation*}
populace. To je \textbf{práh kolektivní imunity}. Pro $R_0 = 3$ je to $\frac{2}{3}$. Spalničky mají $R_0$ zhruba $12$ až $18$, a proto je potřeba proočkovanost přes $92\,\%$. Na obrázku je vidět i druhá stránka věci: vrchol nastane přesně ve chvíli, kdy $S$ klesne na $\frac{1}{R_0}$. Po vrcholu ale epidemie ještě dlouho `dobíhá' a nakazí mnohem víc lidí, než je práh.

\begin{example}
  \textbf{Výška vrcholu.} Vydělíme rovnice pro $I$ a $S$:
  \begin{equation*}
    \frac{\mathrm{d}I}{\mathrm{d}S} = \frac{\beta SI - \gamma I}{-\beta SI} = -1 + \frac{1}{R_0S}\qquad\Longrightarrow\qquad I + S - \frac{1}{R_0}\ln S = \text{konst.}
  \end{equation*}
  Ve vrcholu je $S = \frac{1}{R_0}$, takže
  \begin{equation*}
    I_{\max} = I_0 + S_0 - \frac{1}{R_0} - \frac{1}{R_0}\ln(R_0S_0) = 1 - \frac{1}{3} - \frac{1}{3}\ln 2{,}97\approx 0{,}304.
  \end{equation*}
  Podobně z $\frac{\mathrm{d}S}{\mathrm{d}R} = -R_0S$ plyne $S = S_0\mathrm{e}^{-R_0R}$. Na konci je $I = 0$, $R_\infty = 1 - S_\infty$, a dostaneme rovnici $S_\infty = S_0\mathrm{e}^{-R_0(1 - S_\infty)}$ s řešením $S_\infty\approx 0{,}06$.
\end{example}

\section{Populace a chaos: logistické zobrazení}

Populace hmyzu v jednom roce určuje populaci v roce příštím. Je-li $x_n\in[0, 1]$ velikost populace jako podíl maximální kapacity prostředí, jednoduchý model je
\begin{equation*}
  x_{n+1} = rx_n(1 - x_n).
\end{equation*}
Člen $rx_n$ je rozmnožování, člen $(1 - x_n)$ brzdí růst, když je jedinců moc (nedostatek potravy). Parametr $r\in[0, 4]$ udává plodnost. Robert May v roce 1976 v časopise Nature upozornil, že tento nevinně vypadající vzorec skrývá neuvěřitelně složité chování.

Iterace se dobře kreslí \textbf{pavučinovým diagramem}: z bodu $x_n$ jdeme svisle ke grafu paraboly (to je $x_{n+1}$) a vodorovně k přímce $y = x$ (tím se $x_{n+1}$ stane novým vstupem).

\begin{center}
  \begin{tikzpicture}[scale=4.2]
    \draw[->] (0,0) -- (1.05,0) node[right, font=\small] {$x_n$};
    \draw[->] (0,0) -- (0,1.05) node[above, font=\small] {$x_{n+1}$};
    \draw[gray] (0,0) -- (1,1);
    \draw[thick, domain=0:1, samples=50] plot (\x, {2.8*\x*(1-\x)});
    \draw[colorofwisdom] (0.100,0) -- (0.100,0.252) -- (0.252,0.252) -- (0.252,0.528) -- (0.528,0.528) -- (0.528,0.698) -- (0.698,0.698) -- (0.698,0.590) -- (0.590,0.590) -- (0.590,0.677) -- (0.677,0.677) -- (0.677,0.612) -- (0.612,0.612) -- (0.612,0.665) -- (0.665,0.665) -- (0.665,0.624) -- (0.624,0.624) -- (0.624,0.657) -- (0.657,0.657) -- (0.657,0.631) -- (0.631,0.631) -- (0.631,0.652) -- (0.652,0.652) -- (0.652,0.635) -- (0.635,0.635);
    \node[font=\scriptsize] at (0.5,-0.08) {$r = 2{,}8$: ustálí se v $1 - \frac{1}{r}\approx 0{,}643$};
    \begin{scope}[xshift=1.35cm]
      \draw[->] (0,0) -- (1.05,0) node[right, font=\small] {$x_n$};
      \draw[->] (0,0) -- (0,1.05);
      \draw[gray] (0,0) -- (1,1);
      \draw[thick, domain=0:1, samples=50] plot (\x, {3.8*\x*(1-\x)});
      \draw[colorofwisdom] (0.100,0) -- (0.100,0.342) -- (0.342,0.342) -- (0.342,0.855) -- (0.855,0.855) -- (0.855,0.471) -- (0.471,0.471) -- (0.471,0.947) -- (0.947,0.947) -- (0.947,0.192) -- (0.192,0.192) -- (0.192,0.589) -- (0.589,0.589) -- (0.589,0.920) -- (0.920,0.920) -- (0.920,0.279) -- (0.279,0.279) -- (0.279,0.764) -- (0.764,0.764) -- (0.764,0.684) -- (0.684,0.684) -- (0.684,0.821) -- (0.821,0.821) -- (0.821,0.558) -- (0.558,0.558);
      \node[font=\scriptsize] at (0.5,-0.08) {$r = 3{,}8$: chaos};
    \end{scope}
  \end{tikzpicture}
\end{center}

\begin{example}
  \textbf{Pevný bod a jeho stabilita.} Pevný bod splňuje $x = rx(1 - x)$, tedy $x = 0$ nebo $x^* = 1 - \frac{1}{r}$. Pevný bod iterace $x_{n+1} = f(x_n)$ je stabilní (blízké body se k němu přibližují), pokud $|f'(x^*)| < 1$. Odchylka $\varepsilon$ se totiž podle Taylorova rozvoje změní na $f(x^* + \varepsilon)\approx x^* + f'(x^*)\varepsilon$, tedy se každý krok vynásobí číslem $f'(x^*)$. Zde $f'(x) = r(1 - 2x)$ a
  \begin{equation*}
    f'(x^*) = r\left(1 - 2 + \frac{2}{r}\right) = 2 - r.
  \end{equation*}
  Podmínka $|2 - r| < 1$ dává $1 < r < 3$. Pro $r = 2{,}8$ je $f'(x^*) = -0{,}8$: odchylka se zmenšuje a střídá znaménko, proto pavučina obíhá kolem pevného bodu dovnitř.
\end{example}

Co se děje pro $r > 3$? Pevný bod ztratí stabilitu a populace začne střídat dvě hodnoty: rok hojnosti, rok bídy. Při $r = 1 + \sqrt{6}\approx 3{,}449$ se i tento cyklus rozdvojí na cyklus délky $4$, pak $8$, $16$, \dots\ Zdvojení přicházejí stále rychleji a kolem $r\approx 3{,}5699$ je jich nekonečně mnoho. Za touto hranicí začíná \textbf{chaos}.

\begin{center}
  \begin{tikzpicture}[x=6.4cm, y=5cm]
    \draw[->] (2.5,0) -- (4.07,0) node[right, font=\small] {$r$};
    \draw[->] (2.5,0) -- (2.5,1.05) node[above, font=\small] {$x$};
    \foreach \r in {2.5,3,3.5,4} \node[below, font=\scriptsize] at (\r,0) {$\r$};
    \foreach \y in {0.5,1} \node[left, font=\scriptsize] at (2.5,\y) {$\y$};
    \fill[colorofwisdom] plot[only marks, mark=*, mark size=0.18pt] coordinates {(2.5000,0.600) (2.5050,0.601) (2.5100,0.602) (2.5150,0.602) (2.5200,0.603) (2.5250,0.604) (2.5300,0.605) (2.5350,0.606) (2.5400,0.606) (2.5450,0.607) (2.5500,0.608) (2.5550,0.609) (2.5600,0.609) (2.5650,0.610) (2.5700,0.611) (2.5750,0.612) (2.5800,0.612) (2.5850,0.613) (2.5900,0.614) (2.5950,0.615) (2.6000,0.615) (2.6050,0.616) (2.6100,0.617) (2.6150,0.618) (2.6200,0.618) (2.6250,0.619) (2.6300,0.620) (2.6350,0.620) (2.6400,0.621) (2.6450,0.622) (2.6500,0.623) (2.6550,0.623) (2.6600,0.624) (2.6650,0.625) (2.6700,0.625) (2.6750,0.626) (2.6800,0.627) (2.6850,0.628) (2.6900,0.628) (2.6950,0.629) (2.7000,0.630) (2.7050,0.630) (2.7100,0.631) (2.7150,0.632) (2.7200,0.632) (2.7250,0.633) (2.7300,0.634) (2.7350,0.634) (2.7400,0.635) (2.7450,0.636) (2.7500,0.636) (2.7550,0.637) (2.7600,0.638) (2.7650,0.638) (2.7700,0.639) (2.7750,0.640) (2.7800,0.640) (2.7850,0.641) (2.7900,0.642) (2.7950,0.642) (2.8000,0.643) (2.8050,0.643) (2.8100,0.644) (2.8150,0.645) (2.8200,0.645) (2.8250,0.646) (2.8300,0.647) (2.8350,0.647) (2.8400,0.648) (2.8450,0.649) (2.8500,0.649) (2.8550,0.650) (2.8600,0.650) (2.8650,0.651) (2.8700,0.652) (2.8750,0.652) (2.8800,0.653) (2.8850,0.653) (2.8900,0.654) (2.8950,0.655) (2.9000,0.655) (2.9050,0.656) (2.9100,0.656) (2.9150,0.657) (2.9200,0.658) (2.9250,0.658) (2.9300,0.659) (2.9350,0.659) (2.9400,0.660) (2.9450,0.660) (2.9500,0.661) (2.9550,0.662) (2.9600,0.662) (2.9650,0.663) (2.9700,0.663) (2.9750,0.664) (2.9800,0.664) (2.9800,0.665) (2.9850,0.665) (2.9900,0.664) (2.9900,0.667) (2.9950,0.661) (2.9950,0.662) (2.9950,0.671) (3.0000,0.653) (3.0000,0.654) (3.0000,0.679) (3.0000,0.680) (3.0050,0.642) (3.0050,0.690) (3.0100,0.633) (3.0100,0.699) (3.0150,0.625) (3.0150,0.707) (3.0200,0.619) (3.0200,0.713) (3.0250,0.613) (3.0250,0.718) (3.0300,0.608) (3.0300,0.722) (3.0350,0.603) (3.0350,0.727) (3.0400,0.598) (3.0400,0.731) (3.0450,0.594) (3.0450,0.734) (3.0500,0.590) (3.0500,0.738) (3.0550,0.586) (3.0550,0.741) (3.0600,0.583) (3.0600,0.744) (3.0650,0.579) (3.0650,0.747) (3.0700,0.576) (3.0700,0.750) (3.0750,0.573) (3.0750,0.752) (3.0800,0.570) (3.0800,0.755) (3.0850,0.567) (3.0850,0.758) (3.0900,0.564) (3.0900,0.760) (3.0950,0.561) (3.0950,0.762) (3.1000,0.558) (3.1000,0.765) (3.1050,0.555) (3.1050,0.767) (3.1100,0.553) (3.1100,0.769) (3.1150,0.550) (3.1150,0.771) (3.1200,0.548) (3.1200,0.773) (3.1250,0.545) (3.1250,0.775) (3.1300,0.543) (3.1300,0.777) (3.1350,0.540) (3.1350,0.779) (3.1400,0.538) (3.1400,0.780) (3.1450,0.536) (3.1450,0.782) (3.1500,0.533) (3.1500,0.784) (3.1550,0.531) (3.1550,0.786) (3.1600,0.529) (3.1600,0.787) (3.1650,0.527) (3.1650,0.789) (3.1700,0.525) (3.1700,0.791) (3.1750,0.523) (3.1750,0.792) (3.1800,0.521) (3.1800,0.794) (3.1850,0.519) (3.1850,0.795) (3.1900,0.517) (3.1900,0.797) (3.1950,0.515) (3.1950,0.798) (3.2000,0.513) (3.2000,0.799) (3.2050,0.511) (3.2050,0.801) (3.2100,0.509) (3.2100,0.802) (3.2150,0.507) (3.2150,0.804) (3.2200,0.506) (3.2200,0.805) (3.2250,0.504) (3.2250,0.806) (3.2300,0.502) (3.2300,0.807) (3.2350,0.500) (3.2350,0.809) (3.2400,0.499) (3.2400,0.810) (3.2450,0.497) (3.2450,0.811) (3.2500,0.495) (3.2500,0.812) (3.2550,0.494) (3.2550,0.814) (3.2600,0.492) (3.2600,0.815) (3.2650,0.490) (3.2650,0.816) (3.2700,0.489) (3.2700,0.817) (3.2750,0.487) (3.2750,0.818) (3.2800,0.486) (3.2800,0.819) (3.2850,0.484) (3.2850,0.820) (3.2900,0.482) (3.2900,0.821) (3.2950,0.481) (3.2950,0.823) (3.3000,0.479) (3.3000,0.824) (3.3050,0.478) (3.3050,0.825) (3.3100,0.476) (3.3100,0.826) (3.3150,0.475) (3.3150,0.827) (3.3200,0.474) (3.3200,0.828) (3.3250,0.472) (3.3250,0.829) (3.3300,0.471) (3.3300,0.830) (3.3350,0.469) (3.3350,0.831) (3.3400,0.468) (3.3400,0.832) (3.3450,0.466) (3.3450,0.832) (3.3500,0.465) (3.3500,0.833) (3.3550,0.464) (3.3550,0.834) (3.3600,0.462) (3.3600,0.835) (3.3650,0.461) (3.3650,0.836) (3.3700,0.460) (3.3700,0.837) (3.3750,0.458) (3.3750,0.838) (3.3800,0.457) (3.3800,0.839) (3.3850,0.456) (3.3850,0.840) (3.3900,0.455) (3.3900,0.840) (3.3950,0.453) (3.3950,0.841) (3.4000,0.452) (3.4000,0.842) (3.4050,0.451) (3.4050,0.843) (3.4100,0.449) (3.4100,0.844) (3.4150,0.448) (3.4150,0.845) (3.4200,0.447) (3.4200,0.845) (3.4250,0.446) (3.4250,0.846) (3.4300,0.445) (3.4300,0.847) (3.4350,0.443) (3.4350,0.848) (3.4400,0.442) (3.4400,0.848) (3.4400,0.849) (3.4450,0.440) (3.4450,0.441) (3.4450,0.442) (3.4450,0.849) (3.4500,0.432) (3.4500,0.448) (3.4500,0.847) (3.4500,0.853) (3.4550,0.421) (3.4550,0.460) (3.4550,0.842) (3.4550,0.858) (3.4600,0.413) (3.4600,0.467) (3.4600,0.839) (3.4600,0.861) (3.4650,0.408) (3.4650,0.473) (3.4650,0.837) (3.4650,0.864) (3.4700,0.403) (3.4700,0.479) (3.4700,0.835) (3.4700,0.866) (3.4750,0.399) (3.4750,0.483) (3.4750,0.833) (3.4750,0.868) (3.4800,0.395) (3.4800,0.487) (3.4800,0.832) (3.4800,0.869) (3.4850,0.392) (3.4850,0.491) (3.4850,0.830) (3.4850,0.871) (3.4900,0.389) (3.4900,0.494) (3.4900,0.829) (3.4900,0.872) (3.4950,0.386) (3.4950,0.498) (3.4950,0.828) (3.4950,0.874) (3.5000,0.383) (3.5000,0.501) (3.5000,0.827) (3.5000,0.875) (3.5050,0.380) (3.5050,0.504) (3.5050,0.826) (3.5050,0.876) (3.5100,0.378) (3.5100,0.507) (3.5100,0.825) (3.5100,0.877) (3.5150,0.375) (3.5150,0.509) (3.5150,0.824) (3.5150,0.878) (3.5200,0.373) (3.5200,0.512) (3.5200,0.823) (3.5200,0.879) (3.5250,0.371) (3.5250,0.515) (3.5250,0.823) (3.5250,0.880) (3.5300,0.369) (3.5300,0.517) (3.5300,0.822) (3.5300,0.881) (3.5350,0.367) (3.5350,0.519) (3.5350,0.821) (3.5350,0.882) (3.5400,0.365) (3.5400,0.522) (3.5400,0.820) (3.5400,0.883) (3.5450,0.360) (3.5450,0.366) (3.5450,0.517) (3.5450,0.530) (3.5450,0.817) (3.5450,0.823) (3.5450,0.883) (3.5450,0.885) (3.5500,0.355) (3.5500,0.370) (3.5500,0.506) (3.5500,0.540) (3.5500,0.813) (3.5500,0.828) (3.5500,0.882) (3.5500,0.887) (3.5550,0.351) (3.5550,0.372) (3.5550,0.500) (3.5550,0.546) (3.5550,0.810) (3.5550,0.831) (3.5550,0.881) (3.5550,0.889) (3.5600,0.349) (3.5600,0.374) (3.5600,0.494) (3.5600,0.551) (3.5600,0.809) (3.5600,0.833) (3.5600,0.881) (3.5600,0.890) (3.5650,0.346) (3.5650,0.347) (3.5650,0.372) (3.5650,0.377) (3.5650,0.486) (3.5650,0.495) (3.5650,0.552) (3.5650,0.557) (3.5650,0.806) (3.5650,0.808) (3.5650,0.833) (3.5650,0.837) (3.5650,0.880) (3.5650,0.881) (3.5650,0.891) (3.5700,0.343) (3.5700,0.346) (3.5700,0.365) (3.5700,0.371) (3.5700,0.378) (3.5700,0.481) (3.5700,0.497) (3.5700,0.510) (3.5700,0.554) (3.5700,0.560) (3.5700,0.562) (3.5700,0.804) (3.5700,0.805) (3.5700,0.808) (3.5700,0.827) (3.5700,0.833) (3.5700,0.840) (3.5700,0.880) (3.5700,0.882) (3.5700,0.884) (3.5700,0.891) (3.5700,0.892) (3.5750,0.340) (3.5750,0.348) (3.5750,0.351) (3.5750,0.356) (3.5750,0.363) (3.5750,0.384) (3.5750,0.466) (3.5750,0.513) (3.5750,0.529) (3.5750,0.541) (3.5750,0.548) (3.5750,0.567) (3.5750,0.802) (3.5750,0.811) (3.5750,0.814) (3.5750,0.820) (3.5750,0.826) (3.5750,0.846) (3.5750,0.878) (3.5750,0.885) (3.5750,0.888) (3.5750,0.890) (3.5750,0.891) (3.5750,0.894) (3.5800,0.340) (3.5800,0.347) (3.5800,0.350) (3.5800,0.353) (3.5800,0.359) (3.5800,0.379) (3.5800,0.475) (3.5800,0.519) (3.5800,0.533) (3.5800,0.541) (3.5800,0.548) (3.5800,0.565) (3.5800,0.803) (3.5800,0.811) (3.5800,0.814) (3.5800,0.818) (3.5800,0.824) (3.5800,0.843) (3.5800,0.880) (3.5800,0.887) (3.5800,0.889) (3.5800,0.890) (3.5800,0.891) (3.5800,0.894) (3.5850,0.334) (3.5850,0.344) (3.5850,0.363) (3.5850,0.369) (3.5850,0.395) (3.5850,0.440) (3.5850,0.494) (3.5850,0.508) (3.5850,0.554) (3.5850,0.579) (3.5850,0.580) (3.5850,0.797) (3.5850,0.798) (3.5850,0.809) (3.5850,0.829) (3.5850,0.835) (3.5850,0.857) (3.5850,0.874) (3.5850,0.883) (3.5850,0.886) (3.5850,0.893) (3.5850,0.896) (3.5900,0.336) (3.5900,0.338) (3.5900,0.345) (3.5900,0.355) (3.5900,0.376) (3.5900,0.383) (3.5900,0.462) (3.5900,0.476) (3.5900,0.524) (3.5900,0.550) (3.5900,0.568) (3.5900,0.573) (3.5900,0.801) (3.5900,0.803) (3.5900,0.811) (3.5900,0.822) (3.5900,0.842) (3.5900,0.848) (3.5900,0.879) (3.5900,0.881) (3.5900,0.889) (3.5900,0.892) (3.5900,0.895) (3.5900,0.896) (3.5950,0.338) (3.5950,0.345) (3.5950,0.351) (3.5950,0.370) (3.5950,0.384) (3.5950,0.400) (3.5950,0.425) (3.5950,0.458) (3.5950,0.533) (3.5950,0.548) (3.5950,0.565) (3.5950,0.586) (3.5950,0.795) (3.5950,0.805) (3.5950,0.812) (3.5950,0.819) (3.5950,0.850) (3.5950,0.863) (3.5950,0.872) (3.5950,0.879) (3.5950,0.884) (3.5950,0.890) (3.5950,0.893) (3.5950,0.895) (3.6000,0.332) (3.6000,0.334) (3.6000,0.335) (3.6000,0.376) (3.6000,0.378) (3.6000,0.387) (3.6000,0.449) (3.6000,0.467) (3.6000,0.472) (3.6000,0.572) (3.6000,0.573) (3.6000,0.579) (3.6000,0.799) (3.6000,0.801) (3.6000,0.802) (3.6000,0.845) (3.6000,0.847) (3.6000,0.854) (3.6000,0.878) (3.6000,0.881) (3.6000,0.896) (3.6000,0.897) (3.6050,0.334) (3.6050,0.342) (3.6050,0.345) (3.6050,0.348) (3.6050,0.373) (3.6050,0.384) (3.6050,0.452) (3.6050,0.477) (3.6050,0.536) (3.6050,0.545) (3.6050,0.552) (3.6050,0.571) (3.6050,0.802) (3.6050,0.811) (3.6050,0.814) (3.6050,0.818) (3.6050,0.843) (3.6050,0.853) (3.6050,0.879) (3.6050,0.883) (3.6050,0.892) (3.6050,0.893) (3.6050,0.894) (3.6050,0.897) (3.6100,0.318) (3.6100,0.320) (3.6100,0.331) (3.6100,0.348) (3.6100,0.364) (3.6100,0.433) (3.6100,0.446) (3.6100,0.496) (3.6100,0.515) (3.6100,0.536) (3.6100,0.608) (3.6100,0.614) (3.6100,0.783) (3.6100,0.786) (3.6100,0.819) (3.6100,0.827) (3.6100,0.835) (3.6100,0.856) (3.6100,0.861) (3.6100,0.886) (3.6100,0.892) (3.6100,0.898) (3.6100,0.902) (3.6150,0.329) (3.6150,0.330) (3.6150,0.346) (3.6150,0.385) (3.6150,0.445) (3.6150,0.537) (3.6150,0.538) (3.6150,0.581) (3.6150,0.583) (3.6150,0.614) (3.6150,0.783) (3.6150,0.798) (3.6150,0.799) (3.6150,0.818) (3.6150,0.819) (3.6150,0.856) (3.6150,0.879) (3.6150,0.893) (3.6150,0.899) (3.6200,0.313) (3.6200,0.323) (3.6200,0.356) (3.6200,0.383) (3.6200,0.398) (3.6200,0.408) (3.6200,0.417) (3.6200,0.448) (3.6200,0.467) (3.6200,0.511) (3.6200,0.597) (3.6200,0.626) (3.6200,0.778) (3.6200,0.791) (3.6200,0.830) (3.6200,0.848) (3.6200,0.855) (3.6200,0.867) (3.6200,0.871) (3.6200,0.874) (3.6200,0.880) (3.6200,0.889) (3.6200,0.901) (3.6200,0.905) (3.6250,0.315) (3.6250,0.316) (3.6250,0.335) (3.6250,0.349) (3.6250,0.350) (3.6250,0.439) (3.6250,0.449) (3.6250,0.525) (3.6250,0.528) (3.6250,0.562) (3.6250,0.614) (3.6250,0.619) (3.6250,0.782) (3.6250,0.784) (3.6250,0.808) (3.6250,0.823) (3.6250,0.824) (3.6250,0.855) (3.6250,0.859) (3.6250,0.892) (3.6250,0.897) (3.6250,0.903) (3.6250,0.904) (3.6300,0.305) (3.6300,0.506) (3.6300,0.644) (3.6300,0.770) (3.6300,0.833) (3.6300,0.907) (3.6350,0.302) (3.6350,0.304) (3.6350,0.306) (3.6350,0.485) (3.6350,0.495) (3.6350,0.509) (3.6350,0.520) (3.6350,0.640) (3.6350,0.646) (3.6350,0.650) (3.6350,0.652) (3.6350,0.766) (3.6350,0.767) (3.6350,0.769) (3.6350,0.772) (3.6350,0.827) (3.6350,0.832) (3.6350,0.837) (3.6350,0.842) (3.6350,0.907) (3.6350,0.908) (3.6350,0.909) (3.6400,0.311) (3.6400,0.318) (3.6400,0.344) (3.6400,0.375) (3.6400,0.385) (3.6400,0.414) (3.6400,0.434) (3.6400,0.456) (3.6400,0.457) (3.6400,0.534) (3.6400,0.606) (3.6400,0.626) (3.6400,0.780) (3.6400,0.789) (3.6400,0.821) (3.6400,0.853) (3.6400,0.861) (3.6400,0.869) (3.6400,0.880) (3.6400,0.883) (3.6400,0.894) (3.6400,0.903) (3.6400,0.906) (3.6450,0.295) (3.6450,0.300) (3.6450,0.303) (3.6450,0.343) (3.6450,0.503) (3.6450,0.522) (3.6450,0.528) (3.6450,0.567) (3.6450,0.645) (3.6450,0.652) (3.6450,0.654) (3.6450,0.669) (3.6450,0.758) (3.6450,0.765) (3.6450,0.767) (3.6450,0.770) (3.6450,0.807) (3.6450,0.824) (3.6450,0.827) (3.6450,0.835) (3.6450,0.895) (3.6450,0.908) (3.6450,0.910) (3.6450,0.911) (3.6500,0.293) (3.6500,0.298) (3.6500,0.304) (3.6500,0.360) (3.6500,0.365) (3.6500,0.475) (3.6500,0.490) (3.6500,0.535) (3.6500,0.580) (3.6500,0.641) (3.6500,0.658) (3.6500,0.674) (3.6500,0.755) (3.6500,0.764) (3.6500,0.773) (3.6500,0.802) (3.6500,0.822) (3.6500,0.840) (3.6500,0.846) (3.6500,0.887) (3.6500,0.889) (3.6500,0.908) (3.6500,0.910) (3.6500,0.912) (3.6550,0.289) (3.6550,0.343) (3.6550,0.359) (3.6550,0.384) (3.6550,0.394) (3.6550,0.397) (3.6550,0.400) (3.6550,0.405) (3.6550,0.429) (3.6550,0.489) (3.6550,0.602) (3.6550,0.682) (3.6550,0.752) (3.6550,0.792) (3.6550,0.841) (3.6550,0.864) (3.6550,0.873) (3.6550,0.875) (3.6550,0.876) (3.6550,0.877) (3.6550,0.881) (3.6550,0.890) (3.6550,0.895) (3.6550,0.913) (3.6600,0.293) (3.6600,0.294) (3.6600,0.327) (3.6600,0.343) (3.6600,0.365) (3.6600,0.471) (3.6600,0.528) (3.6600,0.562) (3.6600,0.565) (3.6600,0.574) (3.6600,0.669) (3.6600,0.670) (3.6600,0.759) (3.6600,0.805) (3.6600,0.809) (3.6600,0.810) (3.6600,0.825) (3.6600,0.848) (3.6600,0.888) (3.6600,0.895) (3.6600,0.901) (3.6600,0.912) (3.6650,0.286) (3.6650,0.321) (3.6650,0.322) (3.6650,0.361) (3.6650,0.368) (3.6650,0.440) (3.6650,0.461) (3.6650,0.480) (3.6650,0.586) (3.6650,0.590) (3.6650,0.623) (3.6650,0.691) (3.6650,0.748) (3.6650,0.783) (3.6650,0.798) (3.6650,0.800) (3.6650,0.845) (3.6650,0.853) (3.6650,0.860) (3.6650,0.887) (3.6650,0.889) (3.6650,0.903) (3.6650,0.915) (3.6700,0.281) (3.6700,0.291) (3.6700,0.298) (3.6700,0.338) (3.6700,0.516) (3.6700,0.535) (3.6700,0.538) (3.6700,0.574) (3.6700,0.654) (3.6700,0.660) (3.6700,0.674) (3.6700,0.704) (3.6700,0.741) (3.6700,0.757) (3.6700,0.765) (3.6700,0.768) (3.6700,0.806) (3.6700,0.821) (3.6700,0.823) (3.6700,0.831) (3.6700,0.897) (3.6700,0.911) (3.6700,0.913) (3.6700,0.917) (3.6750,0.277) (3.6750,0.303) (3.6750,0.358) (3.6750,0.371) (3.6750,0.389) (3.6750,0.406) (3.6750,0.450) (3.6750,0.483) (3.6750,0.603) (3.6750,0.639) (3.6750,0.685) (3.6750,0.713) (3.6750,0.737) (3.6750,0.752) (3.6750,0.776) (3.6750,0.793) (3.6750,0.844) (3.6750,0.857) (3.6750,0.874) (3.6750,0.880) (3.6750,0.886) (3.6750,0.891) (3.6750,0.909) (3.6750,0.918) (3.6800,0.272) (3.6800,0.305) (3.6800,0.510) (3.6800,0.555) (3.6800,0.632) (3.6800,0.653) (3.6800,0.669) (3.6800,0.702) (3.6800,0.708) (3.6800,0.721) (3.6800,0.726) (3.6800,0.727) (3.6800,0.729) (3.6800,0.730) (3.6800,0.733) (3.6800,0.740) (3.6800,0.761) (3.6800,0.769) (3.6800,0.780) (3.6800,0.815) (3.6800,0.834) (3.6800,0.856) (3.6800,0.909) (3.6800,0.920) (3.6850,0.268) (3.6850,0.347) (3.6850,0.362) (3.6850,0.467) (3.6850,0.508) (3.6850,0.585) (3.6850,0.593) (3.6850,0.680) (3.6850,0.682) (3.6850,0.712) (3.6850,0.713) (3.6850,0.723) (3.6850,0.738) (3.6850,0.755) (3.6850,0.756) (3.6850,0.799) (3.6850,0.802) (3.6850,0.835) (3.6850,0.851) (3.6850,0.889) (3.6850,0.895) (3.6850,0.917) (3.6850,0.921) (3.6900,0.276) (3.6900,0.283) (3.6900,0.361) (3.6900,0.364) (3.6900,0.437) (3.6900,0.459) (3.6900,0.467) (3.6900,0.595) (3.6900,0.627) (3.6900,0.684) (3.6900,0.695) (3.6900,0.713) (3.6900,0.738) (3.6900,0.749) (3.6900,0.754) (3.6900,0.783) (3.6900,0.798) (3.6900,0.851) (3.6900,0.854) (3.6900,0.863) (3.6900,0.889) (3.6900,0.890) (3.6900,0.916) (3.6900,0.918) (3.6950,0.262) (3.6950,0.265) (3.6950,0.344) (3.6950,0.356) (3.6950,0.479) (3.6950,0.492) (3.6950,0.511) (3.6950,0.593) (3.6950,0.649) (3.6950,0.683) (3.6950,0.702) (3.6950,0.714) (3.6950,0.720) (3.6950,0.745) (3.6950,0.755) (3.6950,0.773) (3.6950,0.799) (3.6950,0.834) (3.6950,0.842) (3.6950,0.847) (3.6950,0.892) (3.6950,0.922) (3.6950,0.923) (3.6950,0.924) (3.7000,0.258) (3.7000,0.277) (3.7000,0.280) (3.7000,0.293) (3.7000,0.489) (3.7000,0.542) (3.7000,0.556) (3.7000,0.649) (3.7000,0.667) (3.7000,0.672) (3.7000,0.702) (3.7000,0.709) (3.7000,0.710) (3.7000,0.741) (3.7000,0.745) (3.7000,0.762) (3.7000,0.764) (3.7000,0.773) (3.7000,0.816) (3.7000,0.822) (3.7000,0.843) (3.7000,0.913) (3.7000,0.918) (3.7000,0.925) (3.7050,0.254) (3.7050,0.256) (3.7050,0.258) (3.7050,0.480) (3.7050,0.509) (3.7050,0.515) (3.7050,0.646) (3.7050,0.656) (3.7050,0.659) (3.7050,0.668) (3.7050,0.702) (3.7050,0.705) (3.7050,0.709) (3.7050,0.764) (3.7050,0.769) (3.7050,0.770) (3.7050,0.775) (3.7050,0.822) (3.7050,0.833) (3.7050,0.836) (3.7050,0.847) (3.7050,0.925) (3.7050,0.926) (3.7100,0.273) (3.7100,0.290) (3.7100,0.312) (3.7100,0.367) (3.7100,0.441) (3.7100,0.545) (3.7100,0.574) (3.7100,0.602) (3.7100,0.669) (3.7100,0.679) (3.7100,0.700) (3.7100,0.720) (3.7100,0.736) (3.7100,0.748) (3.7100,0.759) (3.7100,0.764) (3.7100,0.796) (3.7100,0.809) (3.7100,0.821) (3.7100,0.862) (3.7100,0.889) (3.7100,0.907) (3.7100,0.915) (3.7100,0.920) (3.7150,0.257) (3.7150,0.263) (3.7150,0.356) (3.7150,0.401) (3.7150,0.450) (3.7150,0.469) (3.7150,0.538) (3.7150,0.618) (3.7150,0.637) (3.7150,0.668) (3.7150,0.700) (3.7150,0.710) (3.7150,0.720) (3.7150,0.748) (3.7150,0.765) (3.7150,0.780) (3.7150,0.824) (3.7150,0.852) (3.7150,0.859) (3.7150,0.877) (3.7150,0.893) (3.7150,0.919) (3.7150,0.923) (3.7150,0.925) (3.7200,0.243) (3.7200,0.244) (3.7200,0.325) (3.7200,0.340) (3.7200,0.347) (3.7200,0.408) (3.7200,0.494) (3.7200,0.514) (3.7200,0.558) (3.7200,0.585) (3.7200,0.596) (3.7200,0.684) (3.7200,0.687) (3.7200,0.800) (3.7200,0.805) (3.7200,0.816) (3.7200,0.834) (3.7200,0.842) (3.7200,0.896) (3.7200,0.898) (3.7200,0.903) (3.7200,0.918) (3.7200,0.929) (3.7200,0.930) (3.7250,0.239) (3.7250,0.253) (3.7250,0.284) (3.7250,0.332) (3.7250,0.469) (3.7250,0.503) (3.7250,0.535) (3.7250,0.562) (3.7250,0.590) (3.7250,0.646) (3.7250,0.677) (3.7250,0.686) (3.7250,0.704) (3.7250,0.757) (3.7250,0.777) (3.7250,0.803) (3.7250,0.815) (3.7250,0.826) (3.7250,0.839) (3.7250,0.852) (3.7250,0.901) (3.7250,0.917) (3.7250,0.927) (3.7250,0.931) (3.7300,0.235) (3.7300,0.240) (3.7300,0.256) (3.7300,0.292) (3.7300,0.480) (3.7300,0.505) (3.7300,0.542) (3.7300,0.570) (3.7300,0.650) (3.7300,0.659) (3.7300,0.671) (3.7300,0.680) (3.7300,0.710) (3.7300,0.769) (3.7300,0.771) (3.7300,0.775) (3.7300,0.812) (3.7300,0.824) (3.7300,0.838) (3.7300,0.848) (3.7300,0.914) (3.7300,0.926) (3.7300,0.931) (3.7300,0.932) (3.7350,0.235) (3.7350,0.263) (3.7350,0.360) (3.7350,0.394) (3.7350,0.448) (3.7350,0.481) (3.7350,0.544) (3.7350,0.620) (3.7350,0.651) (3.7350,0.672) (3.7350,0.696) (3.7350,0.706) (3.7350,0.724) (3.7350,0.747) (3.7350,0.775) (3.7350,0.790) (3.7350,0.823) (3.7350,0.848) (3.7350,0.860) (3.7350,0.880) (3.7350,0.892) (3.7350,0.924) (3.7350,0.926) (3.7350,0.932) (3.7400,0.227) (3.7400,0.496) (3.7400,0.657) (3.7400,0.843) (3.7400,0.935) (3.7450,0.304) (3.7450,0.315) (3.7450,0.357) (3.7450,0.378) (3.7450,0.382) (3.7450,0.383) (3.7450,0.384) (3.7450,0.393) (3.7450,0.452) (3.7450,0.582) (3.7450,0.588) (3.7450,0.618) (3.7450,0.792) (3.7450,0.808) (3.7450,0.860) (3.7450,0.881) (3.7450,0.884) (3.7450,0.885) (3.7450,0.886) (3.7450,0.893) (3.7450,0.907) (3.7450,0.911) (3.7450,0.928) (3.7500,0.240) (3.7500,0.268) (3.7500,0.294) (3.7500,0.459) (3.7500,0.563) (3.7500,0.577) (3.7500,0.579) (3.7500,0.646) (3.7500,0.684) (3.7500,0.685) (3.7500,0.718) (3.7500,0.728) (3.7500,0.736) (3.7500,0.742) (3.7500,0.760) (3.7500,0.779) (3.7500,0.809) (3.7500,0.810) (3.7500,0.816) (3.7500,0.857) (3.7500,0.914) (3.7500,0.915) (3.7500,0.922) (3.7500,0.931) (3.7550,0.216) (3.7550,0.279) (3.7550,0.282) (3.7550,0.339) (3.7550,0.426) (3.7550,0.428) (3.7550,0.500) (3.7550,0.602) (3.7550,0.636) (3.7550,0.683) (3.7550,0.693) (3.7550,0.756) (3.7550,0.761) (3.7550,0.799) (3.7550,0.813) (3.7550,0.842) (3.7550,0.869) (3.7550,0.870) (3.7550,0.900) (3.7550,0.918) (3.7550,0.919) (3.7550,0.939) (3.7600,0.249) (3.7600,0.250) (3.7600,0.302) (3.7600,0.321) (3.7600,0.376) (3.7600,0.414) (3.7600,0.445) (3.7600,0.556) (3.7600,0.618) (3.7600,0.632) (3.7600,0.643) (3.7600,0.702) (3.7600,0.706) (3.7600,0.781) (3.7600,0.786) (3.7600,0.793) (3.7600,0.820) (3.7600,0.863) (3.7600,0.874) (3.7600,0.887) (3.7600,0.906) (3.7600,0.912) (3.7600,0.928) (3.7600,0.929) (3.7650,0.212) (3.7650,0.231) (3.7650,0.251) (3.7650,0.319) (3.7650,0.404) (3.7650,0.441) (3.7650,0.458) (3.7650,0.518) (3.7650,0.561) (3.7650,0.630) (3.7650,0.648) (3.7650,0.668) (3.7650,0.708) (3.7650,0.779) (3.7650,0.818) (3.7650,0.835) (3.7650,0.858) (3.7650,0.864) (3.7650,0.878) (3.7650,0.907) (3.7650,0.927) (3.7650,0.928) (3.7650,0.934) (3.7650,0.940) (3.7700,0.205) (3.7700,0.206) (3.7700,0.250) (3.7700,0.288) (3.7700,0.357) (3.7700,0.364) (3.7700,0.417) (3.7700,0.439) (3.7700,0.496) (3.7700,0.613) (3.7700,0.616) (3.7700,0.662) (3.7700,0.707) (3.7700,0.773) (3.7700,0.781) (3.7700,0.844) (3.7700,0.865) (3.7700,0.873) (3.7700,0.892) (3.7700,0.894) (3.7700,0.917) (3.7700,0.929) (3.7700,0.942) (3.7750,0.245) (3.7750,0.246) (3.7750,0.326) (3.7750,0.356) (3.7750,0.370) (3.7750,0.398) (3.7750,0.440) (3.7750,0.535) (3.7750,0.614) (3.7750,0.620) (3.7750,0.698) (3.7750,0.700) (3.7750,0.793) (3.7750,0.796) (3.7750,0.829) (3.7750,0.865) (3.7750,0.880) (3.7750,0.890) (3.7750,0.895) (3.7750,0.905) (3.7750,0.930) (3.7800,0.295) (3.7800,0.297) (3.7800,0.328) (3.7800,0.365) (3.7800,0.375) (3.7800,0.377) (3.7800,0.381) (3.7800,0.396) (3.7800,0.410) (3.7800,0.630) (3.7800,0.635) (3.7800,0.786) (3.7800,0.789) (3.7800,0.834) (3.7800,0.876) (3.7800,0.881) (3.7800,0.886) (3.7800,0.888) (3.7800,0.892) (3.7800,0.904) (3.7800,0.914) (3.7800,0.915) (3.7850,0.239) (3.7850,0.264) (3.7850,0.372) (3.7850,0.378) (3.7850,0.388) (3.7850,0.439) (3.7850,0.576) (3.7850,0.624) (3.7850,0.688) (3.7850,0.702) (3.7850,0.725) (3.7850,0.733) (3.7850,0.735) (3.7850,0.738) (3.7850,0.742) (3.7850,0.754) (3.7850,0.792) (3.7850,0.813) (3.7850,0.866) (3.7850,0.884) (3.7850,0.888) (3.7850,0.890) (3.7850,0.925) (3.7850,0.932) (3.7900,0.189) (3.7900,0.197) (3.7900,0.201) (3.7900,0.224) (3.7900,0.269) (3.7900,0.336) (3.7900,0.475) (3.7900,0.496) (3.7900,0.552) (3.7900,0.580) (3.7900,0.610) (3.7900,0.658) (3.7900,0.681) (3.7900,0.719) (3.7900,0.745) (3.7900,0.765) (3.7900,0.823) (3.7900,0.845) (3.7900,0.853) (3.7900,0.902) (3.7900,0.923) (3.7900,0.937) (3.7900,0.945) (3.7900,0.947) (3.7950,0.205) (3.7950,0.206) (3.7950,0.233) (3.7950,0.353) (3.7950,0.362) (3.7950,0.412) (3.7950,0.438) (3.7950,0.540) (3.7950,0.541) (3.7950,0.618) (3.7950,0.621) (3.7950,0.678) (3.7950,0.679) (3.7950,0.719) (3.7950,0.767) (3.7950,0.828) (3.7950,0.867) (3.7950,0.876) (3.7950,0.893) (3.7950,0.896) (3.7950,0.934) (3.7950,0.942) (3.7950,0.943) (3.8000,0.191) (3.8000,0.200) (3.8000,0.221) (3.8000,0.235) (3.8000,0.278) (3.8000,0.324) (3.8000,0.461) (3.8000,0.529) (3.8000,0.556) (3.8000,0.588) (3.8000,0.608) (3.8000,0.655) (3.8000,0.684) (3.8000,0.689) (3.8000,0.762) (3.8000,0.814) (3.8000,0.822) (3.8000,0.833) (3.8000,0.859) (3.8000,0.906) (3.8000,0.921) (3.8000,0.938) (3.8000,0.944) (3.8000,0.947) (3.8050,0.180) (3.8050,0.191) (3.8050,0.203) (3.8050,0.225) (3.8050,0.274) (3.8050,0.337) (3.8050,0.343) (3.8050,0.467) (3.8050,0.484) (3.8050,0.561) (3.8050,0.588) (3.8050,0.614) (3.8050,0.616) (3.8050,0.701) (3.8050,0.756) (3.8050,0.798) (3.8050,0.851) (3.8050,0.857) (3.8050,0.900) (3.8050,0.902) (3.8050,0.922) (3.8050,0.937) (3.8050,0.947) (3.8050,0.950) (3.8100,0.173) (3.8100,0.181) (3.8100,0.200) (3.8100,0.228) (3.8100,0.231) (3.8100,0.320) (3.8100,0.508) (3.8100,0.526) (3.8100,0.538) (3.8100,0.546) (3.8100,0.566) (3.8100,0.609) (3.8100,0.671) (3.8100,0.676) (3.8100,0.830) (3.8100,0.834) (3.8100,0.841) (3.8100,0.907) (3.8100,0.935) (3.8100,0.936) (3.8100,0.945) (3.8100,0.947) (3.8100,0.950) (3.8100,0.952) (3.8150,0.183) (3.8150,0.205) (3.8150,0.234) (3.8150,0.245) (3.8150,0.355) (3.8150,0.365) (3.8150,0.374) (3.8150,0.423) (3.8150,0.553) (3.8150,0.571) (3.8150,0.623) (3.8150,0.629) (3.8150,0.684) (3.8150,0.706) (3.8150,0.792) (3.8150,0.824) (3.8150,0.873) (3.8150,0.884) (3.8150,0.890) (3.8150,0.893) (3.8150,0.896) (3.8150,0.931) (3.8150,0.934) (3.8150,0.943) (3.8200,0.175) (3.8200,0.184) (3.8200,0.198) (3.8200,0.223) (3.8200,0.236) (3.8200,0.472) (3.8200,0.539) (3.8200,0.551) (3.8200,0.567) (3.8200,0.574) (3.8200,0.606) (3.8200,0.661) (3.8200,0.681) (3.8200,0.689) (3.8200,0.768) (3.8200,0.819) (3.8200,0.830) (3.8200,0.856) (3.8200,0.912) (3.8200,0.934) (3.8200,0.938) (3.8200,0.945) (3.8200,0.949) (3.8200,0.952) (3.8250,0.172) (3.8250,0.185) (3.8250,0.239) (3.8250,0.261) (3.8250,0.544) (3.8250,0.578) (3.8250,0.589) (3.8250,0.696) (3.8250,0.710) (3.8250,0.730) (3.8250,0.736) (3.8250,0.738) (3.8250,0.739) (3.8250,0.740) (3.8250,0.743) (3.8250,0.754) (3.8250,0.788) (3.8250,0.810) (3.8250,0.926) (3.8250,0.933) (3.8250,0.949) (3.8300,0.156) (3.8300,0.505) (3.8300,0.957) (3.8350,0.152) (3.8350,0.495) (3.8350,0.959) (3.8400,0.149) (3.8400,0.488) (3.8400,0.959) (3.8450,0.143) (3.8450,0.154) (3.8450,0.472) (3.8450,0.501) (3.8450,0.958) (3.8450,0.961) (3.8500,0.139) (3.8500,0.141) (3.8500,0.143) (3.8500,0.149) (3.8500,0.156) (3.8500,0.159) (3.8500,0.160) (3.8500,0.461) (3.8500,0.462) (3.8500,0.465) (3.8500,0.472) (3.8500,0.489) (3.8500,0.506) (3.8500,0.514) (3.8500,0.518) (3.8500,0.957) (3.8500,0.958) (3.8500,0.960) (3.8500,0.961) (3.8500,0.962) (3.8550,0.135) (3.8550,0.140) (3.8550,0.153) (3.8550,0.160) (3.8550,0.170) (3.8550,0.449) (3.8550,0.463) (3.8550,0.499) (3.8550,0.519) (3.8550,0.543) (3.8550,0.954) (3.8550,0.956) (3.8550,0.957) (3.8550,0.959) (3.8550,0.962) (3.8550,0.964) (3.8600,0.212) (3.8600,0.241) (3.8600,0.270) (3.8600,0.271) (3.8600,0.300) (3.8600,0.397) (3.8600,0.591) (3.8600,0.603) (3.8600,0.614) (3.8600,0.645) (3.8600,0.700) (3.8600,0.703) (3.8600,0.706) (3.8600,0.761) (3.8600,0.762) (3.8600,0.802) (3.8600,0.806) (3.8600,0.811) (3.8600,0.883) (3.8600,0.915) (3.8600,0.924) (3.8600,0.933) (3.8650,0.129) (3.8650,0.134) (3.8650,0.135) (3.8650,0.158) (3.8650,0.162) (3.8650,0.185) (3.8650,0.218) (3.8650,0.434) (3.8650,0.442) (3.8650,0.449) (3.8650,0.452) (3.8650,0.515) (3.8650,0.526) (3.8650,0.582) (3.8650,0.659) (3.8650,0.868) (3.8650,0.940) (3.8650,0.950) (3.8650,0.953) (3.8650,0.956) (3.8650,0.957) (3.8650,0.964) (3.8650,0.965) (3.8700,0.129) (3.8700,0.133) (3.8700,0.161) (3.8700,0.163) (3.8700,0.179) (3.8700,0.435) (3.8700,0.445) (3.8700,0.447) (3.8700,0.523) (3.8700,0.528) (3.8700,0.529) (3.8700,0.569) (3.8700,0.684) (3.8700,0.725) (3.8700,0.771) (3.8700,0.837) (3.8700,0.949) (3.8700,0.951) (3.8700,0.956) (3.8700,0.964) (3.8700,0.965) (3.8750,0.118) (3.8750,0.123) (3.8750,0.135) (3.8750,0.149) (3.8750,0.211) (3.8750,0.242) (3.8750,0.329) (3.8750,0.404) (3.8750,0.417) (3.8750,0.452) (3.8750,0.480) (3.8750,0.492) (3.8750,0.627) (3.8750,0.644) (3.8750,0.711) (3.8750,0.797) (3.8750,0.855) (3.8750,0.906) (3.8750,0.933) (3.8750,0.942) (3.8750,0.960) (3.8750,0.964) (3.8750,0.967) (3.8750,0.968) (3.8800,0.125) (3.8800,0.190) (3.8800,0.243) (3.8800,0.315) (3.8800,0.357) (3.8800,0.377) (3.8800,0.402) (3.8800,0.425) (3.8800,0.530) (3.8800,0.598) (3.8800,0.637) (3.8800,0.714) (3.8800,0.793) (3.8800,0.837) (3.8800,0.883) (3.8800,0.891) (3.8800,0.897) (3.8800,0.911) (3.8800,0.933) (3.8800,0.948) (3.8800,0.967) (3.8850,0.112) (3.8850,0.120) (3.8850,0.180) (3.8850,0.184) (3.8850,0.203) (3.8850,0.216) (3.8850,0.326) (3.8850,0.386) (3.8850,0.412) (3.8850,0.428) (3.8850,0.485) (3.8850,0.574) (3.8850,0.583) (3.8850,0.628) (3.8850,0.658) (3.8850,0.854) (3.8850,0.874) (3.8850,0.907) (3.8850,0.921) (3.8850,0.941) (3.8850,0.945) (3.8850,0.950) (3.8850,0.951) (3.8850,0.970) (3.8900,0.104) (3.8900,0.124) (3.8900,0.190) (3.8900,0.258) (3.8900,0.351) (3.8900,0.363) (3.8900,0.394) (3.8900,0.421) (3.8900,0.496) (3.8900,0.537) (3.8900,0.688) (3.8900,0.728) (3.8900,0.739) (3.8900,0.745) (3.8900,0.751) (3.8900,0.770) (3.8900,0.835) (3.8900,0.850) (3.8900,0.886) (3.8900,0.900) (3.8900,0.929) (3.8900,0.948) (3.8900,0.967) (3.8900,0.972) (3.8950,0.183) (3.8950,0.193) (3.8950,0.257) (3.8950,0.278) (3.8950,0.423) (3.8950,0.449) (3.8950,0.582) (3.8950,0.607) (3.8950,0.614) (3.8950,0.665) (3.8950,0.709) (3.8950,0.734) (3.8950,0.741) (3.8950,0.745) (3.8950,0.748) (3.8950,0.761) (3.8950,0.781) (3.8950,0.804) (3.8950,0.867) (3.8950,0.923) (3.8950,0.929) (3.8950,0.948) (3.8950,0.951) (3.8950,0.964) (3.9000,0.187) (3.9000,0.194) (3.9000,0.199) (3.9000,0.214) (3.9000,0.262) (3.9000,0.298) (3.9000,0.414) (3.9000,0.584) (3.9000,0.592) (3.9000,0.610) (3.9000,0.622) (3.9000,0.657) (3.9000,0.670) (3.9000,0.724) (3.9000,0.754) (3.9000,0.780) (3.9000,0.817) (3.9000,0.862) (3.9000,0.879) (3.9000,0.916) (3.9000,0.928) (3.9000,0.942) (3.9000,0.946) (3.9000,0.947) (3.9050,0.093) (3.9050,0.105) (3.9050,0.116) (3.9050,0.226) (3.9050,0.327) (3.9050,0.328) (3.9050,0.331) (3.9050,0.368) (3.9050,0.401) (3.9050,0.458) (3.9050,0.468) (3.9050,0.486) (3.9050,0.512) (3.9050,0.683) (3.9050,0.845) (3.9050,0.861) (3.9050,0.864) (3.9050,0.908) (3.9050,0.938) (3.9050,0.969) (3.9050,0.972) (3.9050,0.975) (3.9050,0.976) (3.9100,0.092) (3.9100,0.100) (3.9100,0.193) (3.9100,0.287) (3.9100,0.297) (3.9100,0.314) (3.9100,0.327) (3.9100,0.352) (3.9100,0.379) (3.9100,0.469) (3.9100,0.521) (3.9100,0.587) (3.9100,0.624) (3.9100,0.801) (3.9100,0.816) (3.9100,0.842) (3.9100,0.861) (3.9100,0.891) (3.9100,0.912) (3.9100,0.917) (3.9100,0.920) (3.9100,0.948) (3.9100,0.974) (3.9100,0.976) (3.9150,0.084) (3.9150,0.103) (3.9150,0.137) (3.9150,0.186) (3.9150,0.201) (3.9150,0.303) (3.9150,0.310) (3.9150,0.341) (3.9150,0.361) (3.9150,0.414) (3.9150,0.462) (3.9150,0.532) (3.9150,0.562) (3.9150,0.592) (3.9150,0.629) (3.9150,0.826) (3.9150,0.838) (3.9150,0.880) (3.9150,0.904) (3.9150,0.913) (3.9150,0.946) (3.9150,0.950) (3.9150,0.964) (3.9150,0.973) (3.9200,0.077) (3.9200,0.151) (3.9200,0.205) (3.9200,0.225) (3.9200,0.279) (3.9200,0.294) (3.9200,0.336) (3.9200,0.398) (3.9200,0.428) (3.9200,0.503) (3.9200,0.595) (3.9200,0.638) (3.9200,0.655) (3.9200,0.683) (3.9200,0.788) (3.9200,0.813) (3.9200,0.875) (3.9200,0.885) (3.9200,0.905) (3.9200,0.918) (3.9200,0.939) (3.9200,0.945) (3.9200,0.960) (3.9200,0.980) (3.9250,0.093) (3.9250,0.110) (3.9250,0.261) (3.9250,0.330) (3.9250,0.384) (3.9250,0.408) (3.9250,0.449) (3.9250,0.537) (3.9250,0.659) (3.9250,0.692) (3.9250,0.723) (3.9250,0.731) (3.9250,0.741) (3.9250,0.747) (3.9250,0.752) (3.9250,0.757) (3.9250,0.771) (3.9250,0.787) (3.9250,0.836) (3.9250,0.868) (3.9250,0.882) (3.9250,0.928) (3.9250,0.971) (3.9250,0.976) (3.9300,0.068) (3.9300,0.073) (3.9300,0.085) (3.9300,0.104) (3.9300,0.138) (3.9300,0.266) (3.9300,0.304) (3.9300,0.317) (3.9300,0.366) (3.9300,0.466) (3.9300,0.499) (3.9300,0.550) (3.9300,0.569) (3.9300,0.701) (3.9300,0.768) (3.9300,0.824) (3.9300,0.832) (3.9300,0.851) (3.9300,0.912) (3.9300,0.964) (3.9300,0.973) (3.9300,0.978) (3.9300,0.981) (3.9300,0.982) (3.9350,0.065) (3.9350,0.076) (3.9350,0.145) (3.9350,0.160) (3.9350,0.239) (3.9350,0.275) (3.9350,0.283) (3.9350,0.425) (3.9350,0.489) (3.9350,0.529) (3.9350,0.582) (3.9350,0.625) (3.9350,0.634) (3.9350,0.665) (3.9350,0.715) (3.9350,0.785) (3.9350,0.798) (3.9350,0.802) (3.9350,0.877) (3.9350,0.922) (3.9350,0.958) (3.9350,0.962) (3.9350,0.980) (3.9350,0.983) (3.9400,0.061) (3.9400,0.079) (3.9400,0.226) (3.9400,0.245) (3.9400,0.259) (3.9400,0.287) (3.9400,0.463) (3.9400,0.486) (3.9400,0.615) (3.9400,0.675) (3.9400,0.681) (3.9400,0.688) (3.9400,0.728) (3.9400,0.729) (3.9400,0.755) (3.9400,0.778) (3.9400,0.780) (3.9400,0.807) (3.9400,0.856) (3.9400,0.864) (3.9400,0.929) (3.9400,0.933) (3.9400,0.979) (3.9400,0.984) (3.9450,0.076) (3.9450,0.190) (3.9450,0.199) (3.9450,0.276) (3.9450,0.291) (3.9450,0.305) (3.9450,0.400) (3.9450,0.538) (3.9450,0.597) (3.9450,0.606) (3.9450,0.630) (3.9450,0.634) (3.9450,0.660) (3.9450,0.788) (3.9450,0.799) (3.9450,0.814) (3.9450,0.837) (3.9450,0.886) (3.9450,0.915) (3.9450,0.920) (3.9450,0.942) (3.9450,0.947) (3.9450,0.949) (3.9450,0.980) (3.9500,0.069) (3.9500,0.082) (3.9500,0.111) (3.9500,0.123) (3.9500,0.132) (3.9500,0.228) (3.9500,0.253) (3.9500,0.299) (3.9500,0.389) (3.9500,0.425) (3.9500,0.453) (3.9500,0.536) (3.9500,0.564) (3.9500,0.695) (3.9500,0.743) (3.9500,0.746) (3.9500,0.749) (3.9500,0.827) (3.9500,0.838) (3.9500,0.939) (3.9500,0.965) (3.9500,0.971) (3.9500,0.979) (3.9500,0.982) (3.9550,0.046) (3.9550,0.051) (3.9550,0.103) (3.9550,0.172) (3.9550,0.191) (3.9550,0.202) (3.9550,0.222) (3.9550,0.298) (3.9550,0.309) (3.9550,0.366) (3.9550,0.490) (3.9550,0.521) (3.9550,0.563) (3.9550,0.611) (3.9550,0.637) (3.9550,0.684) (3.9550,0.844) (3.9550,0.855) (3.9550,0.915) (3.9550,0.918) (3.9550,0.940) (3.9550,0.973) (3.9550,0.987) (3.9550,0.988) (3.9600,0.040) (3.9600,0.079) (3.9600,0.150) (3.9600,0.151) (3.9600,0.152) (3.9600,0.205) (3.9600,0.288) (3.9600,0.334) (3.9600,0.414) (3.9600,0.449) (3.9600,0.506) (3.9600,0.507) (3.9600,0.511) (3.9600,0.606) (3.9600,0.645) (3.9600,0.811) (3.9600,0.881) (3.9600,0.907) (3.9600,0.945) (3.9600,0.961) (3.9600,0.980) (3.9600,0.990) (3.9650,0.078) (3.9650,0.213) (3.9650,0.238) (3.9650,0.284) (3.9650,0.286) (3.9650,0.404) (3.9650,0.553) (3.9650,0.610) (3.9650,0.618) (3.9650,0.632) (3.9650,0.664) (3.9650,0.700) (3.9650,0.719) (3.9650,0.736) (3.9650,0.771) (3.9650,0.801) (3.9650,0.807) (3.9650,0.810) (3.9650,0.832) (3.9650,0.885) (3.9650,0.922) (3.9650,0.936) (3.9650,0.943) (3.9650,0.980) (3.9700,0.035) (3.9700,0.054) (3.9700,0.074) (3.9700,0.134) (3.9700,0.190) (3.9700,0.204) (3.9700,0.214) (3.9700,0.273) (3.9700,0.325) (3.9700,0.396) (3.9700,0.446) (3.9700,0.460) (3.9700,0.612) (3.9700,0.644) (3.9700,0.663) (3.9700,0.669) (3.9700,0.788) (3.9700,0.871) (3.9700,0.888) (3.9700,0.910) (3.9700,0.943) (3.9700,0.949) (3.9700,0.981) (3.9700,0.986) (3.9750,0.043) (3.9750,0.050) (3.9750,0.162) (3.9750,0.190) (3.9750,0.206) (3.9750,0.211) (3.9750,0.320) (3.9750,0.347) (3.9750,0.356) (3.9750,0.389) (3.9750,0.466) (3.9750,0.541) (3.9750,0.612) (3.9750,0.651) (3.9750,0.662) (3.9750,0.864) (3.9750,0.890) (3.9750,0.900) (3.9750,0.903) (3.9750,0.912) (3.9750,0.944) (3.9750,0.945) (3.9750,0.987) (3.9750,0.989) (3.9800,0.052) (3.9800,0.102) (3.9800,0.131) (3.9800,0.196) (3.9800,0.211) (3.9800,0.214) (3.9800,0.284) (3.9800,0.365) (3.9800,0.386) (3.9800,0.414) (3.9800,0.454) (3.9800,0.613) (3.9800,0.628) (3.9800,0.661) (3.9800,0.669) (3.9800,0.810) (3.9800,0.882) (3.9800,0.891) (3.9800,0.923) (3.9800,0.943) (3.9800,0.944) (3.9800,0.966) (3.9800,0.974) (3.9800,0.987) (3.9850,0.023) (3.9850,0.044) (3.9850,0.071) (3.9850,0.084) (3.9850,0.088) (3.9850,0.110) (3.9850,0.169) (3.9850,0.261) (3.9850,0.305) (3.9850,0.321) (3.9850,0.389) (3.9850,0.457) (3.9850,0.522) (3.9850,0.560) (3.9850,0.578) (3.9850,0.708) (3.9850,0.769) (3.9850,0.824) (3.9850,0.845) (3.9850,0.868) (3.9850,0.972) (3.9850,0.982) (3.9850,0.989) (3.9850,0.994) (3.9900,0.013) (3.9900,0.050) (3.9900,0.074) (3.9900,0.098) (3.9900,0.189) (3.9900,0.221) (3.9900,0.272) (3.9900,0.325) (3.9900,0.352) (3.9900,0.381) (3.9900,0.436) (3.9900,0.487) (3.9900,0.611) (3.9900,0.662) (3.9900,0.687) (3.9900,0.790) (3.9900,0.858) (3.9900,0.875) (3.9900,0.893) (3.9900,0.911) (3.9900,0.941) (3.9900,0.949) (3.9900,0.981) (3.9900,0.997) (3.9950,0.005) (3.9950,0.020) (3.9950,0.022) (3.9950,0.023) (3.9950,0.078) (3.9950,0.084) (3.9950,0.091) (3.9950,0.131) (3.9950,0.135) (3.9950,0.288) (3.9950,0.309) (3.9950,0.331) (3.9950,0.408) (3.9950,0.454) (3.9950,0.468) (3.9950,0.502) (3.9950,0.590) (3.9950,0.820) (3.9950,0.853) (3.9950,0.885) (3.9950,0.965) (3.9950,0.966) (3.9950,0.995) (3.9950,0.999) (4.0000,0.000)};
    \draw[gray, dashed] (3,0) -- (3,1);
    \draw[gray, dashed] (3.449,0) -- (3.449,1);
    \draw[gray, dashed] (3.5699,0) -- (3.5699,1);
    \node[font=\scriptsize, gray, anchor=south] at (3,1) {$3$};
    \node[font=\scriptsize, gray, anchor=south] at (3.42,1) {$1 + \sqrt{6}$};
    \node[font=\scriptsize, gray, anchor=south] at (3.62,1) {$3{,}57$};
    \node[font=\scriptsize, anchor=west] at (3.79,1.08) {okno periody $3$};
    \draw[->, gray] (3.84,1.05) -- (3.835,0.97);
  \end{tikzpicture}
\end{center}

Na \textbf{bifurkačním diagramu} je pro každé $r$ vyneseno, kterými hodnotami populace prochází po odeznění počátečního přechodu. Vlevo jediná hodnota, pak vidlička na dvě, na čtyři\dots\ a pak `šum'. V chaosu jsou ale vidět i bílá \textbf{okna}, například kolem $r\approx 3{,}83$ se populace náhle zase pravidelně opakuje s periodou $3$. Poměry délek po sobě jdoucích intervalů mezi zdvojeními konvergují k číslu $\delta\approx 4{,}669$, které objevil Mitchell Feigenbaum. Totéž číslo se objeví u každého podobného `hrbolatého' zobrazení, je to univerzální konstanta jako $\pi$.

\begin{example}
  \textbf{Chaos pro $r = 4$ je přesně spočitatelný.} Zapišme $x_n = \sin^2\theta_n$. Pak
  \begin{equation*}
    x_{n+1} = 4\sin^2\theta_n(1 - \sin^2\theta_n) = 4\sin^2\theta_n\cos^2\theta_n = \sin^2(2\theta_n).
  \end{equation*}
  Tedy $\theta_n = 2^n\theta_0$: úhel se každý krok zdvojnásobí. Malá nejistota v $\theta_0$, třeba v desáté desetinné cifře, se za $34$ kroků zvětší $2^{34}\approx 1{,}7\cdot 10^{10}$krát a předpověď je bezcenná. To je \textbf{citlivost na počáteční podmínky}, podstata chaosu. Systém je přitom zcela deterministický, žádná náhoda v něm není.
\end{example}

\section{Slunečnice a zlatý úhel}

Semínka v květu slunečnice vyrůstají jedno po druhém ze středu a každé nové se pootočí o stejný úhel vůči předchozímu. Starší semínka se odsouvají ven. Model (Vogel, 1979): $n$-té semínko leží v polárních souřadnicích
\begin{equation*}
  r_n = c\sqrt{n},\qquad\theta_n = n\cdot\alpha.
\end{equation*}
Odmocnina zajistí, že plocha kruhu s $n$ semínky, $\pi r_n^2 = \pi c^2n$, roste úměrně $n$, takže hustota je všude stejná. Zbývá zvolit úhel $\alpha$. Příroda volí \textbf{zlatý úhel}
\begin{equation*}
  \alpha = 360^\circ\left(1 - \frac{1}{\varphi}\right) = \frac{360^\circ}{\varphi^2}\approx 137{,}508^\circ,\qquad\varphi = \frac{1 + \sqrt{5}}{2}.
\end{equation*}

\begin{center}
  \begin{tikzpicture}
    \foreach \f/\x/\l in {0.381966/0/{$137{,}508^\circ$ (zlatý úhel)}, 0.380556/4.1/{$137^\circ$}, 0.375/8.2/{$135^\circ = \frac{3}{8}$ otáčky}} {
      \begin{scope}[xshift=\x cm]
        \foreach \n in {1,...,300} {
          \pgfmathsetmacro{\t}{360*frac(\n*\f)}
          \fill[colorofwisdom] ({0.105*sqrt(\n)*cos(\t)},{0.105*sqrt(\n)*sin(\t)}) circle (0.05);
        }
        \node[font=\scriptsize] at (0,-2.15) {\l};
      \end{scope}
    }
  \end{tikzpicture}
\end{center}

Se zlatým úhlem (vlevo) vyplní semínka plochu rovnoměrně a bez mezer. Oko v nich vidí spirály, a když je spočítáte, vyjdou sousední Fibonacciho čísla, u skutečných slunečnic typicky $34$ a $55$ nebo $55$ a $89$. Už při $137^\circ$ (uprostřed) se objeví výrazná ramena s mezerami mezi nimi. A při racionálním úhlu $\frac{3}{8}$ otáčky (vpravo) leží každé osmé semínko na stejném paprsku, takže vznikne osm paprsků a obrovské prázdné výseče.

Proč? Pokud je úhel $\alpha$ blízko zlomku $\frac{p}{q}$ otáčky, semínka $n$ a $n + q$ leží skoro na stejném paprsku a vznikne $q$ ramen. Čím lépe lze $\frac{\alpha}{360^\circ}$ aproximovat zlomkem s malým jmenovatelem, tím dřív se ramena objeví a tím víc místa se plýtvá. Ideální úhel je tedy ten, jehož podíl otáčky je \textit{nejhůře aproximovatelný zlomky}, `nejiracionálnější' číslo. Že je to právě $\frac{1}{\varphi^2}$ (a $\varphi$), uvidíme v odbočce.

\section{Hudba: zlomky, které zní dobře}

Tón je kmitání vzduchu, jeho výška je \textbf{frekvence} $f$ v hertzích. Komorní $a^1$ má $440$ Hz. Struna ale nekmitá jen základní frekvencí: zní v ní i \textbf{vyšší harmonické} $2f$, $3f$, $4f$, \dots\ (jejich poměr určuje barvu zvuku, jak jsme viděli u Fourierových řad v kapitole~\ref{kap:rady}). Dva tóny zní \textbf{souzvučně}, když se mnoho jejich harmonických shoduje, a \textbf{nelibě}, když se harmonické míjejí těsně vedle sebe a `drhnou'.

\begin{center}
  \begin{tikzpicture}[x=0.0035cm, y=1cm]
    \draw[->] (0,0) -- (2550,0) node[right, font=\small] {$f$ (Hz)};
    \foreach \f in {0,600,...,2400} \node[below, font=\scriptsize] at (\f,0) {$\f$};
    \foreach \k in {1,...,12} \draw[very thick] ({200*\k},0.25) -- ({200*\k},0.75);
    \foreach \k in {1,...,8} \draw[very thick, colorofwisdom] ({300*\k},1.05) -- ({300*\k},1.55);
    \foreach \k in {1,...,8} \draw[very thick, gray] ({282.84*\k},1.85) -- ({282.84*\k},2.35);
    \foreach \f in {600,1200,1800,2400} {
      \draw[colorofwisdom, dashed] (\f,0.25) -- (\f,1.55);
      \fill[colorofwisdom] (\f,0.9) circle (2pt);
    }
    \node[font=\scriptsize, anchor=east] at (-30,0.5) {$200$ Hz};
    \node[font=\scriptsize, anchor=east] at (-30,1.3) {kvinta $300$ Hz};
    \node[font=\scriptsize, anchor=east] at (-30,2.1) {tritón $200\sqrt{2}$ Hz};
  \end{tikzpicture}
\end{center}

Tón $200$ Hz a tón o \textbf{kvintu} výš, $300$ Hz, s poměrem $3:2$ mají společnou každou třetí, respektive každou druhou harmonickou: $600$, $1200$, $1800$, \dots\ Hz. Ucho slyší shodu a interval zní `čistě'. Ještě čistší je \textbf{oktáva} $2:1$, při ní jsou všechny harmonické horního tónu harmonickými dolního. Tón v poměru $\sqrt{2}$ (tritón, kdysi přezdívaný `ďábel v hudbě') nemá s dolním tónem společnou žádnou harmonickou, protože $\sqrt{2}$ je iracionální. Harmonické se jen míjejí, třeba $1414$ Hz vedle $1400$ Hz, a interval zní napjatě.

\subsection{Rovnoměrné ladění a Pythagorejské koma}

Pythagorejci stavěli stupnici z čistých kvint. Vystoupáme-li ale o $12$ kvint, měli bychom se dostat na stejný tón jako po $7$ oktávách. Nedostaneme:
\begin{equation*}
  \left(\frac{3}{2}\right)^{12} = \frac{531\,441}{4\,096}\approx 129{,}75,\qquad 2^7 = 128,\qquad\frac{3^{12}}{2^{19}} = \frac{531\,441}{524\,288}\approx 1{,}0136.
\end{equation*}
Rozdílu se říká \textbf{Pythagorejské koma}. Nemůže být nula, protože $3^{12}$ je liché a $2^{19}$ sudé. Žádná mocnina čísla $\frac{3}{2}$ není mocninou dvojky, čistými kvintami se kruh nikdy neuzavře.

\begin{definitionbox}
  \textbf{Rovnoměrně temperované ladění} dělí oktávu na $12$ stejných půltónů s poměrem $2^{1/12}\approx 1{,}059\,46$. Tón o $k$ půltónů nad $a^1$ má frekvenci $440\cdot 2^{k/12}$ Hz.
\end{definitionbox}

\begin{center}
  \begin{tabular}{|l|c|c|c|}
    \hline
    interval & půltóny & čistý poměr & $2^{k/12}$\\
    \hline
    oktáva & $12$ & $2$ & $2$\\
    kvinta & $7$ & $1{,}5$ & $1{,}4983$\\
    kvarta & $5$ & $1{,}3333$ & $1{,}3348$\\
    velká tercie & $4$ & $1{,}25$ & $1{,}2599$\\
    \hline
  \end{tabular}
\end{center}

Rovnoměrné ladění je kompromis: všechny intervaly kromě oktávy jsou trochu `falešně', ale ve všech tóninách stejně, takže lze volně modulovat. Kvinta je pokřivená jen o $0{,}1\,\%$, velká tercie ale o $0{,}8\,\%$, což cvičené ucho slyší. Přesný výpočet $2^{1/12}$ publikoval čínský princ Zhu Zaiyu v roce 1584. Proč zrovna $12$? Protože $2^{7/12}$ je nečekaně dobrá aproximace čísla $\frac{3}{2}$. Ekvivalentně je $\frac{7}{12}$ dobrá aproximace čísla $\log_2\frac{3}{2}\approx 0{,}585$, a to je jeden z řetězových zlomků, o nichž bude řeč v odbočce.

\subsection{Zázněje}

Dva tóny s blízkými frekvencemi $f_1$ a $f_2$ se skládají podle vzorce pro součet sinů z kapitoly~\ref{kap:goniometrie}:
\begin{equation*}
  \sin(2\pi f_1t) + \sin(2\pi f_2t) = 2\cos\big(\pi(f_1 - f_2)t\big)\sin\big(\pi(f_1 + f_2)t\big).
\end{equation*}
Slyšíme průměrnou frekvenci $\frac{f_1 + f_2}{2}$, jejíž hlasitost pomalu pulzuje. Obálka $2|\cos(\pi(f_1 - f_2)t)|$ má periodu $\frac{1}{|f_1 - f_2|}$, takže hlasitost kolísá $|f_1 - f_2|$krát za sekundu. Tomu se říká \textbf{zázněje}. Ladiči klavírů je využívají: dokud zázněje slyší, struny nejsou naladěné stejně.

\begin{center}
  \begin{tikzpicture}[x=4.2cm, y=0.75cm]
    \draw[->] (0,0) -- (2.05,0) node[right, font=\small] {$t$ (s)};
    \foreach \t in {0,1,2} \node[below, font=\scriptsize] at (\t,-2) {$\t$};
    \draw[thin, colorofwisdom, domain=0:2, samples=500] plot (\x, {sin(360*8*\x)+sin(360*9*\x)});
    \draw[thick, gray, dashed, domain=0:2, samples=100] plot (\x, {2*cos(180*\x)});
    \draw[thick, gray, dashed, domain=0:2, samples=100] plot (\x, {-2*cos(180*\x)});
    \node[font=\scriptsize, anchor=west] at (2.07,1.4) {obálka $\pm 2\cos(\pi t)$};
    \node[font=\scriptsize, anchor=west] at (2.07,0.6) {$f_1 = 8$, $f_2 = 9$ Hz};
  \end{tikzpicture}
\end{center}

\section{GPS: kružnice, hodiny a relativita}

Každá družice GPS vysílá svou polohu a přesný čas vyslání. Přijímač změří, kdy signál dorazil, a rozdíl časů vynásobený rychlostí světla $c$ je vzdálenost k družici. Bod ve známé vzdálenosti od družice leží na kouli kolem ní, a průsečík několika koulí je naše poloha. Princip se jmenuje \textbf{trilaterace}. Ve 2D jsou to kružnice:

\begin{center}
  \begin{tikzpicture}[scale=0.62]
    \draw[gray!30] (-4,-4) grid (10,8.5);
    \draw[thick, colorofwisdom] (0,0) circle ({sqrt(13)});
    \draw[thick, black] (6,0) circle ({sqrt(13)});
    \draw[thick, gray] (2,5) circle ({sqrt(10)});
    \draw[dashed] (3,-4) -- (3,8.5);
    \draw[dashed] (-2,4) -- (10,-0.8);
    \foreach \p/\l/\pos in {{(0,0)}/A/below left, {(6,0)}/B/below right, {(2,5)}/C/above} {
      \fill \p circle (3pt); \node[\pos, font=\small] at \p {$\l$};
    }
    \fill[colorofwisdom] (3,2) circle (4pt);
    \node[font=\small, anchor=south west] at (3.05,2.05) {$P = [3, 2]$};
    \node[font=\scriptsize, fill=white, inner sep=1pt] at (3,7.8) {$x = 3$};
    \node[font=\scriptsize, fill=white, inner sep=1pt, rotate=-21.8] at (7.5,0.6) {$4x + 10y = 32$};
  \end{tikzpicture}
\end{center}

\begin{example}
  Družice (majáky) jsou v bodech $A = [0, 0]$, $B = [6, 0]$, $C = [2, 5]$ a změřené vzdálenosti jsou $\sqrt{13}$, $\sqrt{13}$, $\sqrt{10}$. Rovnice kružnic
  \begin{equation*}
    x^2 + y^2 = 13,\qquad (x - 6)^2 + y^2 = 13,\qquad (x - 2)^2 + (y - 5)^2 = 10
  \end{equation*}
  jsou kvadratické, ale jejich rozdíly jsou lineární, protože členy $x^2 + y^2$ se odečtou. První minus druhá: $12x - 36 = 0$, tedy $x = 3$. První minus třetí: $4x - 4 + 10y - 25 = 3$, tedy $4x + 10y = 32$ a $y = 2$. Odečtené rovnice jsou přímky z kapitoly~\ref{kap:analyticka} (na obrázku čárkovaně), na nichž leží společné tětivy dvojic kružnic.
\end{example}

Ve 3D potřebujeme tři koule. Proč tedy GPS potřebuje \textit{čtyři} družice? Kvůli hodinám. Družice nesou atomové hodiny, ale přijímač v telefonu má obyčejné křemenné hodiny, které se mohou rozcházet o neznámou odchylku $b$. Chyba hodin o jedinou mikrosekundu je chyba vzdálenosti $c\cdot 10^{-6}\,\mathrm{s}\approx 300$ m. Změřená `pseudovzdálenost' k družici $i$ je proto
\begin{equation*}
  \rho_i = \sqrt{(x - x_i)^2 + (y - y_i)^2 + (z - z_i)^2} + c\cdot b,
\end{equation*}
a neznámé jsou čtyři: $x$, $y$, $z$ a $b$. Na čtyři neznámé potřebujeme čtyři rovnice. Přijímač tak mimochodem zjistí i velmi přesný čas.

\subsection{Relativita v kapse}

Družice GPS obíhají ve výšce asi $20\,200$ km rychlostí asi $3{,}9$ km/s. Na jejich hodinách se projeví oba efekty Einsteinovy relativity:
\begin{itemize}
  \item \textbf{Speciální relativita}: pohybující se hodiny jdou pomaleji, relativně o $\frac{v^2}{2c^2} = \frac{(3\,874)^2}{2\cdot(3\cdot 10^8)^2}\approx 8{,}3\cdot 10^{-11}$. Za den ($86\,400$ s) to dělá asi $-7$ \si{\micro\second}.
  \item \textbf{Obecná relativita}: hodiny v slabším gravitačním poli jdou rychleji, relativně o $\frac{GM}{c^2}\left(\frac{1}{R_Z} - \frac{1}{r}\right)\approx 5{,}3\cdot 10^{-10}$, kde $R_Z = 6\,371$ km je poloměr Země a $r\approx 26\,570$ km poloměr dráhy. Za den asi $+45$ \si{\micro\second}.
\end{itemize}
Celkem jdou hodiny na družici o zhruba $38$ \si{\micro\second} denně rychleji. Vynásobeno rychlostí světla to je asi $11$ km za den. Bez relativistické korekce by se GPS každý den rozladilo zhruba o deset kilometrů. Inženýři proto hodiny na družicích ještě před startem nastavili tak, aby na Zemi tikaly o tento rozdíl pomaleji: místo frekvence $10{,}23$ MHz mají $10{,}229\,999\,995\,43$ MHz. Na oběžné dráze pak jdou správně.

\section{PageRank: náhodný surfař}

Jak poznat, která webová stránka je důležitá? Larry Page a Sergey Brin v roce 1998 navrhli: stránka je důležitá, pokud na ni odkazují důležité stránky. Zní to jako kruh v definici, ale je to přesně rovnice pro vlastní vektor.

\begin{center}
  \begin{tikzpicture}[scale=1.2, pg/.style={draw, thick, circle, minimum size=9mm, font=\small, fill=colorofwisdom!20}]
    \node[pg, minimum size=11mm] (A) at (0,0) {$A$};
    \node[pg, minimum size=8mm] (B) at (2.2,1.5) {$B$};
    \node[pg, minimum size=11mm] (C) at (4.4,0) {$C$};
    \node[pg, minimum size=6mm, fill=white] (D) at (4.4,-2) {$D$};
    \draw[->, thick] (A) -- (B);
    \draw[->, thick] (A) to[bend left=12] (C);
    \draw[->, thick] (B) -- (C);
    \draw[->, thick] (C) to[bend left=12] (A);
    \draw[->, thick] (D) -- (C);
    \node[font=\scriptsize, anchor=east] at (-0.4,0.5) {$0{,}373$};
    \node[font=\scriptsize, anchor=south] at (2.2,1.9) {$0{,}196$};
    \node[font=\scriptsize, anchor=west] at (4.85,0.4) {$0{,}394$};
    \node[font=\scriptsize, anchor=west] at (4.75,-2) {$0{,}0375$};
  \end{tikzpicture}
\end{center}

\begin{definitionbox}
  \textbf{Náhodný surfař} klikne na stránce na náhodný odkaz (každý se stejnou pravděpodobností). Pravděpodobnosti, že je na stránkách $1, \dots, N$, tvoří vektor $\vec{x}$, který se jedním kliknutím změní na $M\vec{x}$, kde $M_{ij} = \frac{1}{\text{počet odkazů ze stránky }j}$, pokud $j$ odkazuje na $i$, jinak $M_{ij} = 0$. Je to \textbf{Markovův řetězec}. \textbf{PageRank} je jeho \textbf{stacionární rozdělení}, vektor s
  \begin{equation*}
    M\vec{x} = \vec{x},\qquad x_1 + \cdots + x_N = 1,
  \end{equation*}
  tedy vlastní vektor matice $M$ s vlastním číslem $1$.
\end{definitionbox}

\begin{example}
  Pro graf na obrázku ($A\to B$, $A\to C$, $B\to C$, $C\to A$, $D\to C$) označme složky vektoru $\vec{x}$ jako $x_A$, $x_B$, $x_C$, $x_D$. Stránka $A$ má dva odkazy, takže každý dostane polovinu její váhy. Stacionární rovnice jsou
  \begin{equation*}
    x_A = x_C,\qquad x_B = \tfrac{1}{2}x_A,\qquad x_C = \tfrac{1}{2}x_A + x_B + x_D,\qquad x_D = 0.
  \end{equation*}
  Na $D$ nic neodkazuje, takže tam surfař nikdy nezůstane. Z normalizace $x_A + \frac{1}{2}x_A + x_A = 1$ dostaneme $\vec{x} = (0{,}4;\ 0{,}2;\ 0{,}4;\ 0)$.
\end{example}

Čistý model má dvě vady. Stránka bez odkazů je past, ze které surfař neuteče, a stránka, na kterou nic neodkazuje, má PageRank nula. Page a Brin proto přidali \textbf{tlumení}: s pravděpodobností $d = 0{,}85$ surfař klikne na odkaz, s pravděpodobností $1 - d = 0{,}15$ se znudí a skočí na úplně náhodnou stránku. Nová matice je
\begin{equation*}
  G = dM + \frac{1 - d}{N}\mathbf{1},
\end{equation*}
kde $\mathbf{1}$ je matice ze samých jedniček. Pro náš graf vyjde $x_D = \frac{0{,}15}{4} = 0{,}0375$ (na $D$ se dá dostat jen skokem) a
\begin{equation*}
  x_A = 0{,}0375 + 0{,}85x_C,\qquad x_B = 0{,}0375 + 0{,}425x_A,\qquad x_C = 0{,}0375 + 0{,}85\left(\tfrac{1}{2}x_A + x_B + 0{,}0375\right).
\end{equation*}
Dosazením dostaneme $x_A = \frac{659}{1769}\approx 0{,}373$, $x_B\approx 0{,}196$, $x_C\approx 0{,}394$ a $x_D = 0{,}0375$ (součet je $1$). Nejdůležitější je $C$, na kterou odkazují tři stránky. Ale $A$ s jediným příchozím odkazem je skoro stejně důležitá, protože ten odkaz přichází od $C$.

\begin{note}
  Proč stacionární vektor existuje a je jediný? Matice $G$ má všechny prvky kladné a součet každého sloupce je $1$. \textbf{Perronova--Frobeniova věta} pak zaručuje, že vlastní číslo $1$ je největší, má jediný kladný vlastní vektor a všechna ostatní vlastní čísla mají absolutní hodnotu nejvýš $d$. Proto \textbf{mocninná metoda} $\vec{x}_{k+1} = G\vec{x}_k$ z libovolného startu konverguje, s chybou klesající zhruba jako $0{,}85^k$. Na celém webu stačí několik desítek iterací. Pro náš graf: $(0{,}25;\ 0{,}25;\ 0{,}25;\ 0{,}25)\to(0{,}25;\ 0{,}144;\ 0{,}569;\ 0{,}038)\to(0{,}521;\ 0{,}144;\ 0{,}298;\ 0{,}038)\to\cdots$
\end{note}

\section{Fraktály: Kochova vločka a délka pobřeží}

\begin{definitionbox}
  \textbf{Kochova vločka} (Helge von Koch, 1904): začneme rovnostranným trojúhelníkem. V každém kroku každou úsečku rozdělíme na třetiny a nad prostřední třetinou postavíme rovnostranný trojúhelník (a jeho základnu odstraníme). Vločka je limita tohoto procesu.
\end{definitionbox}

\begin{center}
  \begin{tikzpicture}[scale=2.5]
    \foreach \k/\x in {0/0, 1/1.2, 2/2.4, 3/3.6} \node[font=\scriptsize] at ({\x+0.5},-0.42) {$n = \k$};
    \draw[thick, fill=colorofwisdom!20] (0.000,0.000) -- (0.500,0.866) -- (1.000,0.000) -- cycle;
    \begin{scope}[shift={(1.2,0)}] \draw[thick, fill=colorofwisdom!20] (0.000,0.000) -- (0.167,0.289) -- (0.000,0.577) -- (0.333,0.577) -- (0.500,0.866) -- (0.667,0.577) -- (1.000,0.577) -- (0.833,0.289) -- (1.000,0.000) -- (0.667,0.000) -- (0.500,-0.289) -- (0.333,0.000) -- cycle; \end{scope}
    \begin{scope}[shift={(2.4,0)}] \draw[thick, fill=colorofwisdom!20] (0.000,0.000) -- (0.056,0.096) -- (0.000,0.192) -- (0.111,0.192) -- (0.167,0.289) -- (0.111,0.385) -- (0.000,0.385) -- (0.056,0.481) -- (0.000,0.577) -- (0.111,0.577) -- (0.167,0.674) -- (0.222,0.577) -- (0.333,0.577) -- (0.389,0.674) -- (0.333,0.770) -- (0.444,0.770) -- (0.500,0.866) -- (0.556,0.770) -- (0.667,0.770) -- (0.611,0.674) -- (0.667,0.577) -- (0.778,0.577) -- (0.833,0.674) -- (0.889,0.577) -- (1.000,0.577) -- (0.944,0.481) -- (1.000,0.385) -- (0.889,0.385) -- (0.833,0.289) -- (0.889,0.192) -- (1.000,0.192) -- (0.944,0.096) -- (1.000,0.000) -- (0.889,0.000) -- (0.833,-0.096) -- (0.778,0.000) -- (0.667,0.000) -- (0.611,-0.096) -- (0.667,-0.192) -- (0.556,-0.192) -- (0.500,-0.289) -- (0.444,-0.192) -- (0.333,-0.192) -- (0.389,-0.096) -- (0.333,0.000) -- (0.222,0.000) -- (0.167,-0.096) -- (0.111,0.000) -- cycle; \end{scope}
    \begin{scope}[shift={(3.6,0)}] \draw[thin, fill=colorofwisdom!20] (0.000,0.000) -- (0.019,0.032) -- (0.000,0.064) -- (0.037,0.064) -- (0.056,0.096) -- (0.037,0.128) -- (0.000,0.128) -- (0.019,0.160) -- (0.000,0.192) -- (0.037,0.192) -- (0.056,0.225) -- (0.074,0.192) -- (0.111,0.192) -- (0.130,0.225) -- (0.111,0.257) -- (0.148,0.257) -- (0.167,0.289) -- (0.148,0.321) -- (0.111,0.321) -- (0.130,0.353) -- (0.111,0.385) -- (0.074,0.385) -- (0.056,0.353) -- (0.037,0.385) -- (0.000,0.385) -- (0.019,0.417) -- (0.000,0.449) -- (0.037,0.449) -- (0.056,0.481) -- (0.037,0.513) -- (0.000,0.513) -- (0.019,0.545) -- (0.000,0.577) -- (0.037,0.577) -- (0.056,0.609) -- (0.074,0.577) -- (0.111,0.577) -- (0.130,0.609) -- (0.111,0.642) -- (0.148,0.642) -- (0.167,0.674) -- (0.185,0.642) -- (0.222,0.642) -- (0.204,0.609) -- (0.222,0.577) -- (0.259,0.577) -- (0.278,0.609) -- (0.296,0.577) -- (0.333,0.577) -- (0.352,0.609) -- (0.333,0.642) -- (0.370,0.642) -- (0.389,0.674) -- (0.370,0.706) -- (0.333,0.706) -- (0.352,0.738) -- (0.333,0.770) -- (0.370,0.770) -- (0.389,0.802) -- (0.407,0.770) -- (0.444,0.770) -- (0.463,0.802) -- (0.444,0.834) -- (0.481,0.834) -- (0.500,0.866) -- (0.519,0.834) -- (0.556,0.834) -- (0.537,0.802) -- (0.556,0.770) -- (0.593,0.770) -- (0.611,0.802) -- (0.630,0.770) -- (0.667,0.770) -- (0.648,0.738) -- (0.667,0.706) -- (0.630,0.706) -- (0.611,0.674) -- (0.630,0.642) -- (0.667,0.642) -- (0.648,0.609) -- (0.667,0.577) -- (0.704,0.577) -- (0.722,0.609) -- (0.741,0.577) -- (0.778,0.577) -- (0.796,0.609) -- (0.778,0.642) -- (0.815,0.642) -- (0.833,0.674) -- (0.852,0.642) -- (0.889,0.642) -- (0.870,0.609) -- (0.889,0.577) -- (0.926,0.577) -- (0.944,0.609) -- (0.963,0.577) -- (1.000,0.577) -- (0.981,0.545) -- (1.000,0.513) -- (0.963,0.513) -- (0.944,0.481) -- (0.963,0.449) -- (1.000,0.449) -- (0.981,0.417) -- (1.000,0.385) -- (0.963,0.385) -- (0.944,0.353) -- (0.926,0.385) -- (0.889,0.385) -- (0.870,0.353) -- (0.889,0.321) -- (0.852,0.321) -- (0.833,0.289) -- (0.852,0.257) -- (0.889,0.257) -- (0.870,0.225) -- (0.889,0.192) -- (0.926,0.192) -- (0.944,0.225) -- (0.963,0.192) -- (1.000,0.192) -- (0.981,0.160) -- (1.000,0.128) -- (0.963,0.128) -- (0.944,0.096) -- (0.963,0.064) -- (1.000,0.064) -- (0.981,0.032) -- (1.000,0.000) -- (0.963,0.000) -- (0.944,-0.032) -- (0.926,0.000) -- (0.889,0.000) -- (0.870,-0.032) -- (0.889,-0.064) -- (0.852,-0.064) -- (0.833,-0.096) -- (0.815,-0.064) -- (0.778,-0.064) -- (0.796,-0.032) -- (0.778,0.000) -- (0.741,0.000) -- (0.722,-0.032) -- (0.704,0.000) -- (0.667,0.000) -- (0.648,-0.032) -- (0.667,-0.064) -- (0.630,-0.064) -- (0.611,-0.096) -- (0.630,-0.128) -- (0.667,-0.128) -- (0.648,-0.160) -- (0.667,-0.192) -- (0.630,-0.192) -- (0.611,-0.225) -- (0.593,-0.192) -- (0.556,-0.192) -- (0.537,-0.225) -- (0.556,-0.257) -- (0.519,-0.257) -- (0.500,-0.289) -- (0.481,-0.257) -- (0.444,-0.257) -- (0.463,-0.225) -- (0.444,-0.192) -- (0.407,-0.192) -- (0.389,-0.225) -- (0.370,-0.192) -- (0.333,-0.192) -- (0.352,-0.160) -- (0.333,-0.128) -- (0.370,-0.128) -- (0.389,-0.096) -- (0.370,-0.064) -- (0.333,-0.064) -- (0.352,-0.032) -- (0.333,0.000) -- (0.296,0.000) -- (0.278,-0.032) -- (0.259,0.000) -- (0.222,0.000) -- (0.204,-0.032) -- (0.222,-0.064) -- (0.185,-0.064) -- (0.167,-0.096) -- (0.148,-0.064) -- (0.111,-0.064) -- (0.130,-0.032) -- (0.111,0.000) -- (0.074,0.000) -- (0.056,-0.032) -- (0.037,0.000) -- cycle; \end{scope}
  \end{tikzpicture}
\end{center}

\begin{example}
  \textbf{Obvod.} Každý krok nahradí každou úsečku čtyřmi úsečkami třetinové délky, takže obvod se vynásobí $\frac{4}{3}$. Při počáteční straně $1$ je obvod po $n$ krocích $3\cdot\left(\frac{4}{3}\right)^n\to\infty$.

  \textbf{Obsah.} Nechť $T$ je obsah počátečního trojúhelníku. V $n$-tém kroku přidáme jeden nový trojúhelník na každou z $3\cdot 4^{n-1}$ úseček, a každý má stranu $3^{-n}$, tedy obsah $\frac{T}{9^n}$. Celkem
  \begin{equation*}
    T + \sum_{n=1}^{\infty}3\cdot 4^{n-1}\cdot\frac{T}{9^n} = T + \frac{T}{3}\sum_{n=1}^{\infty}\left(\frac{4}{9}\right)^{n-1} = T + \frac{T}{3}\cdot\frac{1}{1 - \frac{4}{9}} = T + \frac{3}{5}T = \frac{8}{5}T.
  \end{equation*}
\end{example}

Vločka má tedy \textbf{nekonečný obvod a konečný obsah}. Vejde se do kruhu opsaného počátečnímu trojúhelníku, a přesto je její hranice nekonečně dlouhá. Je to přesně situace, na kterou narazil anglický matematik Lewis Fry Richardson, když zjistil, že délky hranice mezi Španělskem a Portugalskem uváděné v různých zdrojích se výrazně liší. Čím kratší měřítko použijete, tím víc zálivů a výběžků změříte a tím je pobřeží delší. Mandelbrot o tom v roce 1967 napsal slavný článek \textit{Jak dlouhé je pobřeží Británie?} a odpověď zní: záleží na délce pravítka, a bez omezení roste nade všechny meze. \textbf{Paradox pobřeží.}

\subsection{Fraktální dimenze}

Úsečka se skládá z $2$ kopií sebe sama zmenšených na polovinu, čtverec ze $4 = 2^2$ kopií, krychle z $8 = 2^3$ kopií. Obecně: útvar dimenze $D$ se skládá z $N = s^D$ kopií zmenšených $s$krát. Odtud
\begin{equation*}
  D = \frac{\log N}{\log s}.
\end{equation*}

\begin{center}
  \begin{tikzpicture}[scale=1.3]
    \draw[very thick, colorofwisdom] (0,0) -- (1,0);
    \draw[very thick] (1,0) -- (2,0);
    \foreach \x in {0,1,2} \fill (\x,0) circle (1.2pt);
    \node[font=\scriptsize, align=center] at (1,-0.6) {$N = 2$, $s = 2$\\$D = 1$};
    \begin{scope}[xshift=3cm, yshift=-0.4cm]
      \draw[thick, fill=colorofwisdom!25] (0,0) rectangle (0.6,0.6);
      \draw[thick] (0.6,0) rectangle (1.2,0.6) (0,0.6) rectangle (0.6,1.2) (0.6,0.6) rectangle (1.2,1.2);
      \node[font=\scriptsize, align=center] at (0.6,-0.4) {$N = 4$, $s = 2$\\$D = 2$};
    \end{scope}
    \begin{scope}[xshift=5.4cm]
      \draw[very thick, colorofwisdom] (0,0) -- (0.6667,0);
      \draw[very thick] (0.6667,0) -- (1,0.5774) -- (1.3333,0) -- (2,0);
      \foreach \p in {(0,0),(0.6667,0),(1,0.5774),(1.3333,0),(2,0)} \fill \p circle (1.2pt);
      \node[font=\scriptsize, align=center] at (1,-0.6) {$N = 4$, $s = 3$\\$D = \frac{\log 4}{\log 3}\approx 1{,}26$};
    \end{scope}
  \end{tikzpicture}
\end{center}

Kochova křivka se skládá ze $4$ kopií zmenšených $3$krát, takže její dimenze je $\frac{\log 4}{\log 3}\approx 1{,}2619$. Je `víc než čára' (proto má nekonečnou délku), ale `míň než plocha' (proto má nulový obsah). Richardsonova měření dávají pro západní pobřeží Británie dimenzi zhruba $1{,}25$, pro hladké pobřeží Jižní Afriky skoro přesně $1$. Fraktály se objevují všude, kde příroda opakuje stejný proces v různých měřítkách: v rozvětvení plic a cév, v kapradinách, v blescích, v mracích, v říčních sítích.

\section{Včely a šestiúhelníky}

Včela potřebuje rozdělit plochu plástu na buňky stejného obsahu a spotřebovat přitom co nejméně vosku na stěny. Kterým pravidelným mnohoúhelníkem to jde?

Nejdřív: pravidelným $n$-úhelníkem lze rovinu vydláždit jen tehdy, když jeho vnitřní úhel $\frac{(n - 2)\cdot 180^\circ}{n}$ dělí $360^\circ$, tedy když $\frac{2n}{n - 2} = 2 + \frac{4}{n - 2}$ je celé číslo. To nastane, jen když $n - 2$ dělí $4$, tedy pro $n = 3$, $4$, $6$. Jsou jen tři možnosti.

\begin{example}
  Porovnáme obvody pro obsah $1$. Rovnostranný trojúhelník se stranou $a$ má obsah $\frac{\sqrt{3}}{4}a^2$, čtverec $a^2$, pravidelný šestiúhelník (šest rovnostranných trojúhelníků) $\frac{3\sqrt{3}}{2}a^2$. Pro obsah $1$:
  \begin{align*}
    \text{trojúhelník:}\quad & a = \sqrt{4/\sqrt{3}}\approx 1{,}520, & o &= 3a\approx 4{,}559,\\
    \text{čtverec:}\quad & a = 1, & o &= 4,\\
    \text{šestiúhelník:}\quad & a = \sqrt{2/(3\sqrt{3})}\approx 0{,}620, & o &= 6a\approx 3{,}722.
  \end{align*}
  Šestiúhelník vyhrává, o $7\,\%$ proti čtverci a o $18\,\%$ proti trojúhelníku. Kruh by měl obvod $2\sqrt{\pi}\approx 3{,}545$, ale kruhy rovinu nevydláždí.
\end{example}

\begin{center}
  \begin{tikzpicture}[scale=1.15]
    \draw[thick, fill=colorofwisdom!15] (0,0) -- (1.520,0) -- (0.760,1.316) -- cycle;
    \node[font=\scriptsize] at (0.76,-0.35) {$o\approx 4{,}56$};
    \begin{scope}[xshift=2.3cm]
      \draw[thick, fill=colorofwisdom!15] (0,0) rectangle (1,1);
      \node[font=\scriptsize] at (0.5,-0.35) {$o = 4$};
    \end{scope}
    \begin{scope}[xshift=4.5cm, yshift=0.537cm]
      \draw[thick, fill=colorofwisdom!30] (0:0.620) -- (60:0.620) -- (120:0.620) -- (180:0.620) -- (240:0.620) -- (300:0.620) -- cycle;
      \node[font=\scriptsize] at (0,-0.887) {$o\approx 3{,}72$};
    \end{scope}
    \begin{scope}[xshift=7.6cm, yshift=0.537cm]
      \foreach \cx/\cy in {0/0, 0.93/0.537, 0.93/-0.537, -0.93/0.537, -0.93/-0.537, 0/1.074, 0/-1.074} {
        \begin{scope}[shift={(\cx,\cy)}]
          \draw[thick, fill=colorofwisdom!15] (0:0.62) -- (60:0.62) -- (120:0.62) -- (180:0.62) -- (240:0.62) -- (300:0.62) -- cycle;
        \end{scope}
      }
    \end{scope}
    \node[font=\scriptsize] at (2.6,-0.9) {tři pravidelné dlaždice se stejným obsahem $1$ a výřez plástu};
  \end{tikzpicture}
\end{center}

Navíc každá stěna v plástu patří dvěma buňkám, takže na buňku připadá polovina obvodu. Že šestiúhelníky jsou nejlepší i mezi \textit{všemi} rozděleními roviny na stejně velké oblasti, ne jen mezi pravidelnými mnohoúhelníky, tušil už Pappos z Alexandrie ve 4. století. Dokázal to až Thomas Hales v roce 1999.

\section{Odbočka --- Řetězové zlomky a nejiracionálnější číslo}

\begin{definitionbox}
  \textbf{Řetězový zlomek} čísla $x$ vznikne opakováním: odděl celou část, zbytek převrať.
  \begin{equation*}
    x = a_0 + \cfrac{1}{a_1 + \cfrac{1}{a_2 + \cfrac{1}{a_3 + \cdots}}} = [a_0; a_1, a_2, a_3, \dots].
  \end{equation*}
  Useknutím po $k$ krocích dostaneme \textbf{sblížený zlomek} $\frac{p_k}{q_k}$.
\end{definitionbox}

\begin{example}
  $\pi = 3{,}14159\ldots$: celá část $3$, zbytek $0{,}14159$, převrácená hodnota $7{,}0625$. Celá část $7$, zbytek převrácený $15{,}996$, a tak dál: $\pi = [3; 7, 15, 1, 292, \dots]$. Sblížené zlomky jsou $3$, $\frac{22}{7}$, $\frac{333}{106}$, $\frac{355}{113}$, \dots\ Protože další člen $292$ je velký, je $\frac{355}{113} = 3{,}141\,592\,9\ldots$ neobyčejně přesné (chyba $3\cdot 10^{-7}$). Velký člen řetězového zlomku znamená `skvělá aproximace přišla brzo'.
\end{example}

Platí obecná věta: sblížené zlomky splňují $\left|x - \frac{p_k}{q_k}\right| < \frac{1}{q_kq_{k+1}}\leq\frac{1}{a_{k+1}q_k^2}$. Čím větší jsou členy $a_k$, tím lépe se číslo dá aproximovat. Nejhůř aproximovatelné je tedy číslo, jehož členy jsou všechny co nejmenší, tedy rovny $1$:
\begin{equation*}
  x = 1 + \cfrac{1}{1 + \cfrac{1}{1 + \cdots}}\qquad\Longrightarrow\qquad x = 1 + \frac{1}{x}\qquad\Longrightarrow\qquad x^2 - x - 1 = 0\qquad\Longrightarrow\qquad x = \frac{1 + \sqrt{5}}{2} = \varphi.
\end{equation*}
Zlatý řez je $[1; 1, 1, 1, \dots]$. Jeho sblížené zlomky jsou podíly sousedních Fibonacciho čísel:
\begin{equation*}
  \frac{1}{1},\ \frac{2}{1},\ \frac{3}{2},\ \frac{5}{3},\ \frac{8}{5},\ \frac{13}{8},\ \frac{21}{13},\ \dots,\ \frac{F_{k+1}}{F_k}.
\end{equation*}
Opravdu: je-li $\frac{p}{q}$ sblížený zlomek, další je $1 + \frac{1}{p/q} = \frac{p + q}{p}$, což je přesně Fibonacciho pravidlo.

\begin{definitionbox}
  \textbf{Hurwitzova věta} (1891): Pro každé iracionální $x$ existuje nekonečně mnoho zlomků $\frac{p}{q}$ s
  \begin{equation*}
    \left|x - \frac{p}{q}\right| < \frac{1}{\sqrt{5}\,q^2}.
  \end{equation*}
  Konstantu $\sqrt{5}$ nelze zvětšit. Pro $x = \varphi$ by pro větší konstantu existovalo jen konečně mnoho takových zlomků.
\end{definitionbox}

Proč je $\varphi$ hraniční, je vidět z Fibonacciho čísel. Platí Cassiniho identita $F_{k+1}F_{k-1} - F_k^2 = (-1)^k$, odkud $\left|\varphi - \frac{F_{k+1}}{F_k}\right|$ je přibližně $\frac{1}{\sqrt{5}F_k^2}$ a konstantu zlepšit nelze. Zpět ke slunečnici: zlatý úhel je $\frac{1}{\varphi^2} = 2 - \varphi$ otáčky, a i jeho řetězový zlomek $[0; 2, 1, 1, 1, \dots]$ má samé malé členy. Jeho nejlepší aproximace jsou $\frac{1}{3}$, $\frac{2}{5}$, $\frac{3}{8}$, $\frac{5}{13}$, $\frac{8}{21}$, $\frac{13}{34}$, $\frac{21}{55}$, \dots\ Jmenovatele aproximací jsou Fibonacciho čísla, a protože ramena se objevují podle jmenovatelů, spirály v květu jsou počítány Fibonacciho čísly. Žádná aproximace není `příliš dobrá', a tak žádná soustava ramen nezačne dominovat a nenechá mezery.

\section{Příběh --- Pumpa na Broad Street}

V srpnu 1854 udeřila v londýnské čtvrti Soho cholera. Během deseti dnů zemřelo přes pět set lidí v okruhu několika ulic, celkem jich bylo přes šest set. Lékařská věda té doby věřila, že cholera se šíří `miasmatem', zkaženým vzduchem z hnijících odpadků. Lékař John Snow, známý mimo jiné tím, že při porodu prince Leopolda podal královně Viktorii chloroform, si myslel něco jiného: že nemoc se přenáší vodou.

Snow obcházel domy, zapisoval adresy zemřelých a zanášel je jako čárky do mapy čtvrti. Mapa ukázala zřetelné shluky kolem jediné veřejné pumpy na Broad Street. Výjimky jeho teorii jen potvrzovaly. Chudobinec s více než pěti sty obyvateli, jen kousek od pumpy, měl vlastní studnu a téměř žádné oběti. Dělníci z místního pivovaru pili pivo, které dostávali jako součást mzdy, a nezemřel ani jeden. A naopak vdova z Hampsteadu, kilometry daleko, zemřela na choleru, i když v jejím okolí nikdo jiný neonemocněl. Vodu z Broad Street měla tak ráda, že si ji nechávala vozit v lahvích.

Sedmého září předložil Snow svá data místní radě. Rada mu příliš nevěřila, ale druhý den nechala z pumpy odmontovat rukojeť. Epidemie v tu dobu už sama slábla, protože mnoho obyvatel uprchlo. Až později místní farář Henry Whitehead, který chtěl Snowovu teorii původně vyvrátit, dohledal pravděpodobný počátek: pětiměsíční holčičku z domu číslo $40$, jejíž matka prala plenky a vylévala vodu do žumpy jen pár stop od studny. Holčička zemřela.

Úřední zpráva Snowovo vysvětlení odmítla a rukojeť byla po čase vrácena. Snow zemřel v roce 1858, ve čtyřiceti pěti letech, aniž by se dočkal uznání. Teprve po letech se jeho mapa stala jedním z nejslavnějších obrázků v dějinách vědy a Snow zakladatelem epidemiologie. Byla to statistika a geometrie, ne mikroskop, co poprvé ukázalo, jak se nemoc šíří.

\section{Shrnutí a otázky}

\begin{problem}
  Vypočítejte:
  \begin{enumerate}[label=(\alph*)]
    \item Nemoc má $R_0 = 5$. Jaká část populace musí být imunní, aby epidemie nepropukla?
    \item Pro která $r$ je pevný bod logistického zobrazení kladný a stabilní? Jaký je pro $r = 2{,}5$?
    \item Kolikrát se zvětší obvod Kochovy vločky po $10$ krocích? Jaká je fraktální dimenze útvaru, který se skládá z $8$ kopií zmenšených $3$krát?
    \item Kolik hertzů má tón o kvintu (7 půltónů) nad $a^1 = 440$ Hz v rovnoměrném ladění a kolik v čistém ladění? Kolik záznějů za sekundu spolu vydají?
  \end{enumerate}
\end{problem}

\begin{solution}
  \begin{enumerate}[label=(\alph*)]
    \item $1 - \frac{1}{5} = 80\,\%$.
    \item Pro $1 < r < 3$, kdy $|2 - r| < 1$. Pro $r = 2{,}5$ je $x^* = 1 - \frac{1}{2{,}5} = 0{,}6$.
    \item $\left(\frac{4}{3}\right)^{10}\approx 17{,}8$krát. $D = \frac{\log 8}{\log 3}\approx 1{,}89$ (to je Sierpińského koberec).
    \item $440\cdot 2^{7/12}\approx 659{,}26$ Hz a $440\cdot\frac{3}{2} = 660$ Hz. Zázněje: asi $0{,}74$ za sekundu (u třetí harmonické tónu $a^1$ a druhé harmonické kvinty, $1320$ proti $1318{,}5$ Hz, je to $1{,}5$ záznějů za sekundu).
  \end{enumerate}
\end{solution}

\begin{problem}
  Odpovězte vlastními slovy:
  \begin{enumerate}
    \item Odkud se v modelu SIR bere člen $\beta SI$ a co znamená $R_0$?
    \item Odvoďte práh kolektivní imunity $1 - \frac{1}{R_0}$. Proč epidemie nakazí víc lidí, než je tento práh?
    \item Proč je pevný bod logistického zobrazení stabilní jen pro $1 < r < 3$? Co se stane pak?
    \item Co je citlivost na počáteční podmínky a jak ji ukazuje substituce $x = \sin^2\theta$ pro $r = 4$?
    \item Proč slunečnice používá zlatý úhel a proč se v ní objevují Fibonacciho čísla?
    \item Proč zní kvinta souzvučně a proč se kruh čistých kvint nikdy neuzavře?
    \item Proč potřebuje GPS čtyři družice a proč musí počítat s relativitou?
    \item Jak souvisí PageRank s vlastním vektorem matice? K čemu je tlumení $d = 0{,}85$?
    \item Jak může mít Kochova vločka nekonečný obvod a konečný obsah? Co je fraktální dimenze?
    \item Proč včely stavějí šestiúhelníkové buňky?
  \end{enumerate}
\end{problem}
